\subsection{极值点与重心}

命题 3. --- 设 K 是 Hausdorff 局部凸空间 E 的一个紧凸子集，x 是 K 的一个点。K 上每个以 x 为重心、总质量为 1 的正测度 \(\mu\)，都属于以 x 为重心、总质量为 1 的离散正测度所成之集的淡闭包。

设 U 是 \(\mu\) 对淡拓扑的一个邻域；我们可以假定 U 由 K 上满足下列条件的测度 \(\nu\) 组成：

\[
\left| \mu (f _ {i}) - \nu (f _ {i}) \right| \leqslant \delta\tag{1}
\]

其中 \(f_{i} \in \mathcal{C}(\mathrm{K}; \mathbf{C})\) ( \(1 \leqslant i \leqslant p\) ) 是有限个函数，\(\delta > 0\) 是一个数。对每个点 \(a \in K\)，存在 E 中 0 的一个闭凸邻域 \(V_{a}\)，使得

\[
\left| f _ {i} (\mathbf {y}) - f _ {i} (\mathbf {a}) \right| \leqslant \delta / 2\tag{2}
\]

对 \(1 \leqslant i \leqslant p\) 及每个 \(\mathbf{y} \in \mathrm{W}_{\mathbf{a}} = \mathrm{K} \cap (\mathbf{a} + \mathrm{V}_{\mathbf{a}})\) 成立。由于 \(\mathbf{K}\) 是紧的，存在有限个点 \(\mathbf{a}_j\) ( \(1 \leqslant j \leqslant r\) )（属于 \(\mathbf{K}\)），使得诸 \(\mathrm{W}_{\mathbf{a}_j}\) 构成 \(\mathbf{K}\) 的一个覆盖 ( \(1 \leqslant j \leqslant r\) )。取一个连续单位分解 \((g_j)_{1 \leqslant j \leqslant r}\)（在 \(\mathbf{K}\) 上，从属于覆盖 \((\mathrm{W}_{\mathbf{a}_j})\)），并令 \(\alpha_j = \mu(g_j)\)（对一切 \(j\)）；若 \(\alpha_j \neq 0\)，令 \(\mu_j = \alpha_j^{-1} g_j \cdot \mu\)，若 \(\alpha_j = 0\)，令 \(\mu_j = \varepsilon_{\mathbf{a}_j}\)。每个测度 \(\mu_j\) 都是正的、总质量为 1，且其支集含于紧凸集 \(\mathrm{W}_{\mathbf{a}_j}\)；此外，按定义，

\[
\mu = \sum_ {j = 1} ^ {r} \alpha_ {j} \mu_ {j}\tag{3}
\]

因为 \(g_{j} \cdot \mu = 0\)（若 \(\mu(g_{j}) = 0\)）；诸 \(\alpha_{j}\) 均 \(\geqslant 0\)，且

\[
\sum_ {j = 1} ^ {r} \alpha_ {j} = \sum_ {j = 1} ^ {r} \mu (g _ {j}) = \mu \left(\sum_ {j = 1} ^ {r} g _ {j}\right) = \mu (1) = 1.
\]

设 \(\mathbf{x}_j\) 是 \(\mu_j\) 的重心，它属于 \(\mathrm{W}_{\mathbf{a}_j}\)（No.~1, 命题 1），并考虑离散测度 \(\nu = \sum_{j=1}^{r} \alpha_j \varepsilon_{\mathbf{x}_j}\)；它是正的且总质量为 1，其重心是 \(\sum_{j=1}^{r} \alpha_j \mathbf{x}_j\)，由 (3)，后者也是 \(\mu\) 的重心，从而等于 \(\mathbf{x}\)。此外，由 (2)，有 \(|f_i(\mathbf{y}) - f_i(\mathbf{a}_j)| \leqslant \delta/2\)（对一切 \(\mathbf{y} \in \mathrm{W}_{\mathbf{a}_j}\) 及一切 \(i\)）；于是，由于 \(\operatorname{Supp}(\mu_j) \subset \mathrm{W}_{\mathbf{a}_j}\)，有 \(|\mu_j(f_i) - f_i(\mathbf{a}_j)| \leqslant \delta/2\)（对 \(1 \leqslant i \leqslant p\)）。另一方面，由于 \(\mathbf{x}_j \in \mathrm{W}_{\mathbf{a}_j}\)，还有

\[
\left| \varepsilon_ {\mathbf {x} _ {j}} (f _ {i}) - f _ {i} (\mathbf {a} _ {j}) \right| \leqslant \delta / 2
\]

对 \(1 \leqslant i \leqslant p\) 成立；从而 \(|\mu_j(f_i) - \varepsilon_{\mathbf{x}_j}(f_i)| \leqslant \delta\)（对一切 \(i\) 和 \(j\)）。由于诸 \(\alpha_j\) 均 \(\geqslant 0\) 且其和为 1，由 (3) 及 \(\nu\) 的定义推出 \(\nu\) 满足不等式 (1)。

证毕。

推论. --- 设 \(K'\) 是 K 的一个紧子集，且 K 是 \(K'\) 的闭凸包络。为使 \(x \in K'\) 是 K 的极值点，必须且只须 \(\varepsilon_{x}\) 是 \(K'\) 上以 x 为重心、总质量为 1 的唯一正测度。

设 \(\mathbf{x}\) 是 \(K\) 的一个极值点；为证明 \(\varepsilon_{\mathbf{x}}\) 是 \(K'\) 上以 \(\mathbf{x}\) 为重心、总质量为 1 的唯一正测度，由命题 3，只须看到离散测度 \(\nu\)（在 \(K'\) 上，为正的、总质量为 1，且以 \(\mathbf{x}\) 为重心）所成之集化为 \(\varepsilon_{\mathbf{x}}\)。但这样的测度 \(\nu\) 形如 \(\sum_{i=1}^{r} \lambda_i \varepsilon_{\mathbf{x}_i}\)，其中 \(\lambda_i > 0\)（对 \(1 \leqslant i \leqslant r\)）且 \(\sum_{i=1}^{r} \lambda_i = 1\)，而 \(\mathbf{x}\) 是 \(\nu\) 的重心这一假设可以写成 \(\mathbf{x} = \sum_{i=1}^{r} \lambda_i \mathbf{x}_i\)。由于 \(\mathbf{x}\) 是极值点，这就蕴涵 \(\mathbf{x}_i = \mathbf{x}\)（对一切 \(i\)），从而 \(\nu = \varepsilon_{\mathbf{x}}\)。

反之，假定 \(\varepsilon_{x}\) 是 \(K'\) 上以 x 为重心、总质量为 1 的唯一正测度，我们来证明 x 是极值点。若不然，就会存在 K 的两个不同点 \(x', x''\) 以及一个实数 \(\lambda\)，使得 \(0 < \lambda < 1\) 且 \(\mathbf{x} = \lambda\mathbf{x}' + (1 - \lambda)\mathbf{x}''\)。由命题 1，\(x'\)（相应地，\(x''\)）是某个正测度 \(\mu'\)（相应地，\(\mu''\)）的重心，该测度在 \(K'\) 上且总质量为 1。于是 x 是 \(\lambda\mu' + (1 - \lambda)\mu''\) 的重心。因此 \(\lambda\mu' + (1 - \lambda)\mu'' = \varepsilon_{x}\)。从而 \(\mu'\) 与 \(\mu''\) 都与 \(\varepsilon_{x}\) 成比例，于是 \(x' = x'' = x\)，这是荒谬的。

定理 1 (Choquet). --- 设 E 是 R 上的一个 Hausdorff 局部凸空间，K 是 E 的一个可度量化紧凸子集，M 是 K 的极值点所成之集。集 M 是 K 中一族可数开集的交，且 K 的每个点都是 K 上某个测度 \(\mu\) 的重心，其中 \(\mu(K - M) = 0\)。

为证明第一个断言，用 I 表示区间 {[}0,1{]}（\(\mathbf{R}\) 中的）；由于 K 是紧且可度量化的，\(\mathrm{K} \times \mathrm{K} \times \mathrm{I}\) 亦然；\(\mathrm{K} \times \mathrm{K} \times \mathrm{I}\) 中由三元组 \((\mathbf{x}, \mathbf{y}, \lambda)\)（满足 \(\mathbf{x} \neq \mathbf{y}\) 且 \(0 < \lambda < 1\)）组成的子集 U 在 \(\mathrm{K} \times \mathrm{K} \times \mathrm{I}\) 中是开集，因此存在一列闭集 \((\mathrm{F}_n)\)，它们在 \(\mathrm{K} \times \mathrm{K} \times \mathrm{I}\) 中，其并为 U（GT, IX, §2, No.~5, 命题 7）。映射 \(q: \mathrm{K} \times \mathrm{K} \times \mathrm{I} \to \mathrm{K}\)（由 \(q(\mathbf{x}, \mathbf{y}, \lambda) = \lambda \mathbf{x} + (1 - \lambda) \mathbf{y}\) 定义）是连续的，且由极值点的定义，\(\mathrm{K} - \mathrm{M} = q(\mathrm{U}) = \bigcup_{n} q(\mathrm{F}_n)\)；但 \(\mathrm{F}_n\) 是紧的，因为它在 \(\mathrm{K} \times \mathrm{K} \times \mathrm{I}\) 中是闭集，因此 \(q(\mathrm{F}_n)\) 是紧的，从而在 K 中是闭的；于是集 \(\mathrm{U}_n = \mathrm{K} - q(\mathrm{F}_n)\) 在 K 中是开集，且有 \(\mathrm{M} = \bigcap \mathrm{U}_n\)。

在证明的后续部分，我们将用 u 表示 K 上的一个连续凸数值函数，用 \(G \subset E \times R\) 表示 u 的图象，用 S 表示 G 在 \(E \times R\) 中的闭凸包络。

引理 1. --- 设 \(\overline{u}\) 是 E 上在 K 上 \(\geqslant u\) 的连续仿射线性函数的下包络。那么，S 是由点 \((\mathbf{a}, b) \in \mathrm{E} \times \mathbf{R}\)（满足 \(\mathbf{a} \in \mathrm{K}\) 且 \(u(\mathbf{a}) \leqslant b \leqslant \overline{u}(\mathbf{a})\)）组成之集。

由 Hahn-Banach 定理，为使 \((\mathbf{a}, b)\) 属于 S，必须且只须 \(h(\mathbf{a}, b) \geqslant 0\) 对 \(E \times R\) 上每个满足 \(h(\mathbf{x}, u(\mathbf{x})) \geqslant 0\)（对 \(x \in K\)）的连续仿射线性函数 h 成立。令 \(h(\mathbf{x}, t) = f(\mathbf{x}) - \lambda t\)，其中 f 是 E 上的一个连续仿射线性函数，我们看到关系 \((\mathbf{a}, b) \in S\) 等价于下列性质：关系

\[
f (\mathbf {x}) \geqslant \lambda u (\mathbf {x}) \quad \text { for   all } \mathbf {x} \in K\tag{4}
\]

蕴涵

\[
f (\mathbf {a}) \geqslant \lambda b.\tag{5}
\]

先设 \(\lambda = 0\)；由 Hahn-Banach 定理，对每个连续仿射线性函数 \(f\)（\(\mathbf{E}\) 上的），(4) 蕴涵 (5) 这一事实等价于关系 \(\mathbf{a} \in \mathbf{K}\)。其次，设 \(\lambda = -1\) 并把 \(f\) 换成 \(-f\)；那么关系 \(f(\mathbf{x}) \leqslant u(\mathbf{x})\)（在 \(\mathbf{K}\) 中）必须蕴涵 \(f(\mathbf{a}) \leqslant b\)。但由于 \(u\) 在 \(\mathbf{K}\) 上是凸的且连续的，\(u(\mathbf{a})\) 等于 \(f(\mathbf{a})\)（对函数 \(f\)——\(\mathbf{E}\) 上满足 \(f(\mathbf{x}) \leqslant u(\mathbf{x})\)（在 \(\mathbf{K}\) 上）的连续仿射线性函数——而言）的上确界（TVS, II, §5, No.~4, 命题 5）；由此得到关系 \(b \geqslant u(\mathbf{a})\)。最后，设 \(\lambda = 1\)；这时说 (4) 蕴涵 (5)，按定义即表示 \(b \leqslant \overline{u}(\mathbf{a})\)。这就证明了引理，因为关系 (4)（相应地，(5)）等价于两边同乘一个标量 \(>0\) 所得的关系。

引理 2. --- 若 u 在 K 上是严格凸的，则对 K 的每个非极值点 x，有 \(u(\mathbf{x}) < \overline{u}(\mathbf{x})\)。

因为这时存在 K 的两个不同点 y,z，使得 \(\mathbf{x} = (\mathbf{y} + \mathbf{z})/2\)，于是由 u 的严格凸性有 \(u(\mathbf{x}) < (u(\mathbf{y}) + u(\mathbf{z}))/2\)。若 f 是 E 上的一个仿射线性函数，则 \(f(\mathbf{x}) = (f(\mathbf{y}) + f(\mathbf{z}))/2\)；把这一关系应用于 E 上在 K 上 \(\geqslant u\) 的连续仿射线性函数 f，便得

\[
\overline {{{u}}} (\mathbf {x}) \geqslant \left(\overline {{{u}}} (\mathbf {y}) + \overline {{{u}}} (\mathbf {z})\right) / 2 \geqslant \left(u (\mathbf {y}) + u (\mathbf {z})\right) / 2 > u (\mathbf {x}).
\]

这两个引理既已建立，于是（引理 1）有 \((\mathbf{a},\overline{u} (\mathbf{a}))\in S\)（对一切 \(\mathbf{a}\in K\)）。由于 G 是 K 在连续映射 \(\mathbf{x}\mapsto (\mathbf{x},u(\mathbf{x}))\) 下的象，故是紧的，由 No.~1 的命题 1，存在一个正测度 \(\nu\)（在 G 上、总质量为 1、以 \((\mathbf{a},\overline{u} (\mathbf{a}))\) 为重心）。由于投影 \(\mathrm{pr}_1\) 在 G 上的限制是 G 到 K 上的一个同胚，测度 \(\nu\) 可借助这一同胚转移，从而给出 K 上的一个测度 \(\mu\)（正的且总质量为 1），使得

\[
\mathbf {a} = \int \mathbf {x} d \mu (\mathbf {x}) \quad \mathrm{and} \quad \overline {{{u}}} (\mathbf {a}) = \int u (\mathbf {x}) d \mu (\mathbf {x}).\tag{6}
\]

关系 (6) 中的第一个表明 a 是 \(\mu\) 的重心。函数 \(\overline{u}\) 在 K 上是上半连续且有界的，因而是 \(\mu\)-可积的（\(\S4\), No.~4, 命题 5 的推论 1）；此外，由于函数 \(-\overline{u}\) 按定义是凸的，由 No.~1 命题 2 的推论，有

\[
\overline {{{u}}} (\mathbf {a}) \geqslant \int \overline {{{u}}} (\mathbf {x}) d \mu (\mathbf {x}),\tag{7}
\]

于是，与公式 (6) 中的第二个相比较，便得

\[
\int u (\mathbf {x}) d \mu (\mathbf {x}) \geqslant \int \overline {{{u}}} (\mathbf {x}) d \mu (\mathbf {x}).\tag{8}
\]

但由于在 K 上 \(u(\mathbf{x}) \leqslant \overline{u}(\mathbf{x})\)，关系 (8) 蕴涵 \(u(\mathbf{x}) = \overline{u}(\mathbf{x})\) 对 \(\mu\) 而言几乎处处成立。考虑到引理 2 可见，只要证明了下面的引理，定理 1 即告完成：

引理 3. --- 设 E 是 R 上的一个 Hausdorff 局部凸空间，K 是 E 的一个可度量化紧凸子集。那么，K 上存在一个严格凸的数值函数。

因为 Banach 空间 \(\mathcal{C}(\mathrm{K};\mathbf{R})\) 是可分的（GT, X, §3, No.~3, 定理 1），故子空间 \(\mathcal{A}\)（由 E 上连续仿射线性函数在 K 上的限制组成，含于 \(\mathcal{C}(\mathrm{K};\mathbf{R})\)）也是可分的。于是设 \((h_n)\) 是 \(\mathcal{A}\) 中一个稠密序列，并设 \(\alpha_{n} > \sup_{\mathbf{x}\in \mathbb{K}}|h_n(\mathbf{x})|\)。那么每个函数 \(h_n^2 / n^2\alpha_n^2\) 在 K 中都是凸的（TVS, II, §2, No.~8, 例），且以 \(h_n^2 / n^2\alpha_n^2\) 为一般项的级数是正规收敛的，因此其和 \(u\) 在 K 中是连续且凸的。剩下要证 \(u\) 是严格凸的，为此只须证明：对 K 的任意两个不同点 \(\mathbf{x},\mathbf{x}'\)，存在一个整数 \(n\)，使得 \(h_n^2\) 在以 \(\mathbf{x},\mathbf{x}'\) 为端点的线段上的限制是严格凸的；而为此只须 \(h_n(\mathbf{x})\neq h_n(\mathbf{x}')\)（同前）。现在，存在一个函数 \(h\in \mathcal{A}\) 使得 \(h(\mathbf{x})\neq h(\mathbf{x}')\)（TVS, II, §4, No.~1, 命题 2 的推论 1），又由于序列 \((h_n)\) 在 \(\mathcal{A}\) 中稠密，故存在一个 \(n\) 使得 \(h_n(\mathbf{x})\neq h_n(\mathbf{x}')\)。

证毕。

推论. --- 设 E 是 R 上的一个 Hausdorff 局部凸空间，C 是 E 中以 0 为顶点的真凸锥，它对 E 的弱化一致结构所诱导的一致结构是完备且可度量化的。设 M 是 C 的极值母线之并。对每个 \(x \in C\)，存在 C 的一个紧凸子集 K 以及 K 上一个总质量为 1 的测度 \(\lambda \geqslant 0\)，使得 \(K - (M \cap K)\) 是 \(\lambda\)-可忽略的，且 \(\lambda\) 的重心等于 x。

因为 x 属于 C 的某个帽 K（TVS, II, §7, No.~2, 命题 5），且 M∩K 包含 K 的极值点所成之集（同前，命题 4 的推论 1）。于是只须应用定理 1。

\subsection{应用：I. 连续实函数的向量空间}

设 X 是一个非空紧空间，H 是 Banach 空间 \(\mathcal{C}(\mathrm{X};\mathbf{R})\) 的一个线性子空间，它包含常数且分离 X 的点（GT, X, §4, No.~1, 定义 1）。我们给 \(\mathcal{H}\) 赋以由 \(\mathcal{C}(\mathrm{X};\mathbf{R})\) 的拓扑诱导的赋范空间拓扑，并用 \(\mathcal{H}'\) 表示这一赋范空间的对偶。对每个 \(x\in \mathrm{X}\)，映射 \(f\mapsto f(x)\) 是 \(\mathcal{H}\) 上的一个连续线性形式（\(\mathcal{H}\) 上 Dirac 测度 \(\varepsilon_{x}\) 的限制），因而是 \(\mathcal{H}'\) 中一个将记作 \(i_{\mathcal{H}}(x)\) 的元素，使得

\[
\langle f, i _ {\mathcal {H}} (x) \rangle = f (x)\tag{9}
\]

对每个函数 \(f \in \mathcal{H}\) 及每个 \(x \in X\) 成立。

映射 \(i_{\mathcal{H}}\)（从 X 到 \(\mathcal{H}'\) 中）当 \(\mathcal{H}'\) 赋以弱拓扑 \(\sigma(\mathcal{H}', \mathcal{H})\) 时是单射且连续的；第二个断言由定义及 (9) 立即得出；至于第一个，注意若 \(x, x'\) 是 X 的两个不同点，则由假设存在一个函数 \(h \in \mathcal{H}\) 使得 \(h(x) \neq h(x')\)，于是由 (9)，\(\langle h, i_{\mathcal{H}}(x) \rangle \neq \langle h, i_{\mathcal{H}}(x') \rangle\)，从而更有 \(i_{\mathcal{H}}(x) \neq i_{\mathcal{H}}(x')\)。因此象 \(i_{\mathcal{H}}(\mathrm{X})\) 是 \(\mathcal{H}'\) 的一个紧子集（对弱拓扑），且 \(i_{\mathcal{H}}\) 是 X 到 \(i_{\mathcal{H}}(\mathrm{X})\) 上的一个同胚。

命题 4. --- (i) \(i_{\mathcal{H}}(X)\) 在 \(\mathcal{H}'\) 中的闭凸包络 C（对弱拓扑 \(\sigma(\mathcal{H}', \mathcal{H})\)）是紧的。

\begin{enumerate}
\def\labelenumi{(\roman{enumi})}
\setcounter{enumi}{1}
\tightlist
\item
  为使点 \(i_{\mathcal{H}}(x)\) 是 \(C\) 的极值点，必须且只须：唯一的正测度 \(\lambda\)（\(X\) 上的）满足
\end{enumerate}

\[
h (x) = \int h d \lambda\tag{10}
\]

（对每个函数 \(h \in \mathcal{H}\)；这特别蕴涵 \(\lambda\) 的总质量为 1，因为 \(1 \in \mathcal{H}\)）就是 Dirac 测度 \(\varepsilon_x\)。

由此得出，对每个 \(h \in \mathcal{H}\)，函数 \(z' \mapsto \langle h, z' \rangle\)（\(C\) 上的）在 \(C\) 的至少一个极值点处达到其上确界（TVS, II, §7, No.~1, 命题 1），且这个点属于 \(i_{\mathcal{H}}(X)\)（同前，命题 2 的推论）。

\begin{enumerate}
\def\labelenumi{(\roman{enumi})}
\item
  由 (9)，\(\| i_{\mathcal{H}}(x)\| \leqslant 1\) 在赋范空间 \(\mathcal{H}'\) 中成立，换言之 \(i_{\mathcal{H}}(\mathrm{X})\) 是有界的，而这一断言由 \(\mathcal{H}'\) 赋以弱拓扑 \(\sigma (\mathcal{H}',\mathcal{H})\) 时是拟完备的这一事实得出（TVS, III, §4, No.~2, 定理 1 的推论 5）。
\item
  每个质量为 1 的正测度 \(\mu\)（在 \(i_{\mathcal{H}}(X)\) 上）都可借助同胚 \(i_{H}\) 通过结构转移而来自 X 上一个质量为 1 的正测度 \(\lambda\)，其中 Dirac 测度 \(\varepsilon_{i_{\mathcal{H}(x)}}\) 来自 \(\varepsilon_{x}\)。说 \(\mu\) 以 \(i_{\mathcal{H}(x)}\) 为重心，按定义即表示
\end{enumerate}

\[
\int_ {\mathrm{X}} \langle h, i _ {\mathcal {H}} (z) \rangle d \lambda (z) = \langle h, i _ {\mathcal {H}} (x) \rangle
\]

对每个函数 \(h \in H\) 成立。考虑到 (9)，断言 (ii) 无非是 No.~2 命题 3 的推论中关于 \(i_{\mathcal{H}}(x)\) 是 C 的极值点的判别准则的翻译。

我们把满足命题 4 的条件 (ii) 的点 \(x \in \mathbf{X}\) 称为 \(\mathcal{H}\)-极值点；用 \(\mathrm{Ch}_{\mathcal{H}}(\mathbf{X})\)（或简作 \(\mathrm{Ch}(\mathbf{X})\)）表示这些点所成之集，用 \(\check{\mathrm{S}}_{\mathcal{H}}(\mathbf{X})\)（或简作 \(\check{\mathrm{S}}(\mathbf{X})\)）表示 \(\mathrm{Ch}_{\mathcal{H}}(\mathbf{X})\) 在 \(\mathbf{X}\) 中的闭包。

命题 5. --- 每个函数 \(h \in \mathcal{H}\) 都在至少一个 \(\mathcal{H}\)-极值点处达到其上确界。

设 \(x\) 是 \(X\) 的一个点，\(h\) 是 \(\mathcal{H}\) 中的一个函数。关系 \(h(z) \leqslant h(x)\)（对一切 \(z \in X\)）可以写成 \(\langle h, i_{\mathcal{H}}(z) \rangle \leqslant \langle h, i_{\mathcal{H}}(x) \rangle\)（对一切 \(z \in X\)），因而它表示 \(\mathcal{H}'\) 中方程为 \(\langle h, t' \rangle = \langle h, i_{\mathcal{H}}(x) \rangle\) 的弱闭超平面是 \(i_{\mathcal{H}}(X)\) 的一个支撑超平面。已知（TVS, II, §7, No.~1, 命题 1 的推论）这样的超平面至少包含 \(i_{\mathcal{H}}(X)\) 的闭凸包络的一个极值点，而这样的点 \(i_{\mathcal{H}}(y)\) 按定义是一个 \(\mathcal{H}\)-极值点 \(y\) 的象；因此 \(h(y)\) 等于 \(h\) 在 \(X\) 中的上确界。

命题 6. --- 对每个点 \(x \in X\)，下列性质等价：

\begin{enumerate}
\def\labelenumi{\alph{enumi})}
\item
  \(x\) 是 \(\mathcal{H}\)-极值点。
\item
  对 x 在 X 中的每个开邻域 U 及每个 \(\varepsilon > 0\)，存在 H 中的一个函数 \(h \geqslant 0\)，使得 \(h(x) \leqslant \varepsilon\) 且 \(h(y) \geqslant 1\)（对每个 \(y \in X - U\)）。
\end{enumerate}

设 \(x\) 是 \(X\) 的任一点，\(f\) 是 \(\mathcal{C}(X; \mathbf{R})\) 中的一个函数；已知（TVS, II, §3, No.~1, 命题 1）诸数 \(\lambda(f)\)（对 \(X\) 上满足 \(\lambda(h) = h(x)\)——对每个函数 \(h \in \mathcal{H}\)——的一切正测度而取的）的下确界，等于诸数 \(h(x)\) 的上确界，其中 \(h\) 取遍函数 \(h \in \mathcal{H}\)（满足 \(h \leqslant f\)）所成之集。假定 \(x\) 是 \(\mathcal{H}\)-极值点；于是由命题 4, (ii)，对每个函数 \(f \in \mathcal{C}(X; \mathbf{R})\)，

\[
f (x) = \sup _ {h \in \mathscr {H}, h \leqslant f} h (x).\tag{11}
\]

为证明 \(a)\) 蕴涵 \(b)\)，我们取 \(f\) 为 X 到 \([0,1]\) 中的一个连续映射，其支集含于 U，且 \(f(x)=1\)；于是由 (11)，存在一个函数 \(h'\in\mathscr{H}\) 使得 \(h'\leqslant f\) 且 \(h'(x)\geqslant 1-\varepsilon\)。由于 \(1\in\mathscr{H}\)，函数 \(h=1-h'\) 满足 \(b)\) 的条件。

反之，假定条件 \(b)\) 成立；这一条件蕴涵 \(1 - h \leqslant \varphi_{\mathrm{U}}\)；于是对每个正测度 \(\lambda\)（\(X\) 上的，满足条件 (10)），我们有

\[
\lambda (\mathrm{U}) = \lambda (\varphi_ {\mathrm{U}}) \geqslant \lambda (1 - h) = 1 - h (x) \geqslant 1 - \varepsilon .
\]

由于按假设这一关系对每个 \(\varepsilon > 0\) 及 x 的每个开邻域 U 都成立，于是得到

\[
\lambda (\{x \}) = \inf _ {U} \lambda (U) \geqslant 1 - \varepsilon
\]

对每个 \(\varepsilon > 0\) 成立，因此 \(\lambda(\{x\}) = 1\)。由于 \(\lambda\) 是正的且总质量为 1，必有 \(\lambda = \varepsilon_x\)，这就证明了 \(x\) 是 \(\mathcal{H}\)-极值点，依据是命题 4, (ii)。

命题 7. --- 设 F 是 X 的一个闭子集。下列性质等价：

\begin{enumerate}
\def\labelenumi{\alph{enumi})}
\item
  F 包含 \(\check{\mathrm{S}}_{\mathcal{H}}(\mathrm{X})\)
\item
  对每个函数 \(h \in \mathcal{H}\)，集 \(F\) 都与 \(X\) 中使 \(h\) 达到其上确界的点所成之集相交。
\item
  对每个点 \(x \in X\)，存在一个总质量为 1 的正测度 \(\mu\)（在 \(X\) 上），使得 \(\operatorname{Supp}(\mu) \subset F\) 且 \(h(x) = \int h d\mu\)（对每个函数 \(h \in \mathcal{H}\)）。
\end{enumerate}

设 \(G = i_{\mathcal{H}}(F)\)。条件 \(a\)) 表示 \(G\) 包含 \(C\) 的极值点所成之集。条件 \(b\)) 表示 \(G\) 与 \(i_{\mathcal{H}}(X)\) 同 \(i_{\mathcal{H}}(X)\) 的每个闭支撑超平面的交相遇。最后，条件 \(c\)) 表示 \(i_{\mathcal{H}}(X)\) 的每个点都是某个支集含于 \(G\) 的测度的重心；由 No.~1 命题 1，这也等价于说 \(i_{\mathcal{H}}(X)\) 的闭凸包络等于 \(G\) 的闭凸包络。于是条件 \(a\))、\(b\)) 与 \(c\)) 的等价性由 TVS, II, §7, No.~1, 命题 2 的推论得出。

命题 8. --- 假定 X 是可度量化的。那么，X 的 H-极值点所成之集 \(\operatorname{Ch}_{\mathcal{H}}(\mathbf{X})\) 是 X 中一族可数开集的交，且对每个 \(x \in X\)，存在 X 上一个总质量为 1 的正测度 \(\mu\)，使得

\[
\mu (\mathrm{X} - \mathrm{Ch} _ {\mathcal {H}} (\mathrm{X})) = 0 \text { and } \int h d \mu = h (x)
\]

对每个 \(h \in \mathcal{H}\) 成立。

这就是 No.~2 的定理 1 借助同胚 \(x \mapsto i_{\mathcal{H}}(x)\) 通过结构转移得到的翻译，如同命题 5 中那样。

当把 H 换成一个由定义在 X 上、取值于 \(R \cup \{+\infty\}\) 的下半连续函数组成的集 P 时，本节的若干结果可以推广，这里假定 P 包含常数且满足 \(P + P \subset P\)（习题 2）。

*例. --- 取 X 为单位球 \(\| \mathbf{x} \| \leqslant 1\)（\(\mathbf{R}^3\) 中的），并设 \(\mathcal{H}\) 是 X 上连续函数的一个向量空间，它包含 \(\mathbf{R}^3\) 上仿射线性函数在 X 上的限制，且满足“极大值原理”，即对每个非常数函数 \(h \in \mathcal{H}\)，X 中使 \(h\) 达到其上确界的点所成之集含于球面 \(\mathbf{S}_2\)。于是由命题 5 与命题 7 容易得出 \(\mathrm{Ch}_{\mathcal{H}}(\mathrm{X}) = \check{\mathrm{S}}_{\mathcal{H}}(\mathrm{X}) = \mathbf{S}_2\)。满足前述条件的向量空间 \(\mathcal{H}\) 的一个重要例子是在 X 上连续且在开球 \(\| \mathbf{x} \| < 1\) 中调和的函数所成之集。对这些函数，人们证明：质量为 1 的正测度 \(\mu\)（满足 \(\operatorname{Supp}(\mu) \subset \mathbf{S}_2\) 且 \(h(\mathbf{x}) = \int h d\mu\)——对一切 \(h \in \mathcal{H}\)）当 \(\| \mathbf{x} \| < 1\) 时由 Poisson 公式给出

\[
d \mu (\mathbf {z}) = \frac {1 - \| \mathbf {z} \| ^ {2}}{\| \mathbf {z} - \mathbf {x} \| ^ {3}} d \sigma (\mathbf {z}),
\]

其中 \(\sigma\) 是 \(\mathbf{S}_2\) 上在正交群下不变且满足 \(\sigma(\mathbf{S}_2) = 1\) 的测度（Ch. VII, §3, 习题 8）。*

\subsection{应用：II. 连续复函数的向量空间}

设 X 是一个非空紧空间，H 是复 Banach 空间 \(\mathcal{C}(\mathrm{X};\mathbf{C})\) 的一个线性子空间，它包含常数且分离 X 的点。诸 \(\mathcal{R}(f)\)（函数 \(f \in H\) 的实部）所成之集是一个线性子空间 \(H_{r}\)（实向量空间 \(\mathcal{C}(\mathrm{X};\mathbf{R})\) 的）；对每个 \(f \in H\)，集 \(H_{r}\) 还包含 \(\mathcal{I}(f) = \mathcal{R}(-if)\)；由此推出 \(H_{r}\) 分离 X 的点，因为关系 \(h(x) = h(y)\)（对一切 \(h \in H_{r}\)）蕴涵 \(\mathcal{R}(f(x)) = \mathcal{R}(f(y))\) 与 \(\mathcal{I}(f(x)) = \mathcal{I}(f(y))\)，从而 \(f(x) = f(y)\)（对一切 \(f \in H\)）。X 中的 \(H_{r}\)-极值点仍称为 H-极值点，其全体记作 \(\operatorname{Ch}_{\mathcal{H}}(\mathrm{X})\)，后者的闭包记作 \(\check{\operatorname{S}}_{\mathcal{H}}(\mathrm{X})\)。命题 5 与命题 7 的类似物如下：

命题 9. --- 对每个函数 \(f \in \mathcal{H}\)，\(\mathrm{Ch}_{\mathcal{H}}(\mathrm{X})\) 与使 \(|f|\) 达到其上确界的点所成之集相交。

我们可以限于 \(f\) 不是常数 0 的情形。设 \(a\) 是 \(X\) 中使 \(|f|\) 达到其上确界的一个点，并令 \(g = f / f(a)\)；那么 \(g(a) = 1\) 且 \(|g(x)| \leqslant 1\)（对一切 \(x \in X\)），于是

\[
\mathscr {R} (g (a)) = 1 \quad \text { and } \quad \mathscr {R} (g (x)) \leqslant 1 \quad \text { for   all } x \in X.
\]

把 No.~3 的命题 5 应用于 \(\mathcal{H}_r\)，存在 \(b \in \mathrm{Ch}_{\mathcal{H}}(\mathrm{X})\) 使 \(\mathcal{R}(g(x))\) 在该点达到其上确界 1，于是 \(|g(b)| = 1\)（因为 \(|g(b)| \leqslant 1\)）；由此得出 \(|f(b)| = |f(a)| \geqslant |f(x)|\)（对一切 \(x \in \mathrm{X}\)）。

命题 10. --- 设 F 是 X 的一个闭子集。下列性质等价：

\begin{enumerate}
\def\labelenumi{\alph{enumi})}
\item
  F 包含 \(\check{\mathrm{S}}_{\mathcal{H}}(\mathrm{X})\)
\item
  对每个函数 \(f \in \mathcal{H}\)，\(F\) 都与 \(X\) 中使 \(|f|\) 达到其上确界的点所成之集相交。
\item
  对每个点 \(x \in X\)，存在 X 上一个总质量为 1 的正测度 \(\mu\)，使得 \(\operatorname{Supp}(\mu) \subset \operatorname{F}\) 且 \(f(x) = \int f \, d\mu\)（对每个函数 \(f \in H\)）。
\end{enumerate}

我们来证明条件 \(a)\) 与 \(c)\) 的等价性：设 \(f = f_{1} + if_{2}\)，其中 \(f_{1}, f_{2}\) 属于 \(\mathcal{H}_{r}\)；关系 \(f(x) = \int f d\mu\) 等价于两个关系 \(f_{1}(x) = \int f_{1} d\mu\) 与 \(f_{2}(x) = \int f_{2} d\mu\)；于是对 \(\mathcal{H}_{r}\) 只须应用 No.~3 命题 7 中条件 \(a)\) 与 \(c)\) 的等价性。\(a)\) 蕴涵 \(b)\) 这一事实由命题 9 得出。我们来证明 \(b)\) 蕴涵 \(a)\)；这就是要看到：若 \(b)\) 成立，则对每个 \(h \in \mathcal{H}_{r}\)，\(F\) 都与 \(h\) 在 \(X\) 中达到其下确界的点所成之集相交。条件 \(b)\) 蕴涵 \(F\) 非空；由于 \(F\) 是紧的，存在 \(a \in F\) 使得 \(h(a) \leqslant h(y)\)（对一切 \(y \in F\)）。设 \(f \in \mathcal{H}\) 满足 \(h = \mathcal{R}(f)\)；对每个 \(\varepsilon > 0\)，函数 \(g = f - h(a) + \varepsilon\) 属于 \(\mathcal{H}\)，且

\[
\mathcal {R} (g (y)) = h (y) - h (a) + \varepsilon \geqslant \varepsilon
\]

对一切 \(y \in F\) 成立。用 c 表示 \(|g|\) 在 X 中的上确界，并令 \(b = c^{2}/2\varepsilon\)；对每个 \(y \in F\)，

\[
\left| g (y) - b \right| ^ {2} = \left| g (y) \right| ^ {2} - 2 b \mathcal {R} (g (y)) + b ^ {2} \leqslant c ^ {2} - 2 b \varepsilon + b ^ {2} = b ^ {2},
\]

换言之，函数 \(|g - b|\) 在 F 中的上确界 \(\leqslant b\)。由于 \(g - b \in \mathcal{H}\)，关于 F 的假设蕴涵 \(|g - b| \leqslant b\)，于是

\[
b ^ {2} \geqslant | g - b | ^ {2} = | g | ^ {2} - 2 b \mathcal {R} (g) + b ^ {2}
\]

从而 \(\mathcal{R}(g) \geqslant |g|^{2}/2b \geqslant 0\)；由于 \(\mathcal{R}(g) = h - h(a) + \varepsilon\) 且 \(\varepsilon > 0\) 任意，我们有 \(h \geqslant h(a)\)，而 \(h(a)\) 是 h 在 X 中的下确界，证明至此完成。

注. --- 若 f 是一个连续实函数，则使 \(|f|\) 达到其上确界的点也是函数 f 与 -f 之一达到其上确界的点。对满足 No.~3 中假设的连续实函数的向量空间 H 而言，命题 9 与命题 10 因而是命题 5 与命题 7 的平凡推论。

\subsection{应用：III. 连续函数的代数}

引理 4. --- 设 X 是一个紧空间，H 是 Banach 空间 \(\mathcal{C}(\mathrm{X};\mathbf{C})\)（相应地，\(\mathcal{C}(\mathrm{X};\mathbf{R})\)）的一个闭线性子空间。设 a 是 X 的一个具有可数基本邻域系的点；假定对满足 0 < c < d < 1 的任何数 c 和 d 以及 a 的任何开邻域 U，存在一个 \(f \in H\)，使得

\[
| f | \leqslant 1, \quad | f (a) | \geqslant d, \quad | f (x) | \leqslant c \text {   for   all   } x \in X - U.\tag{12}
\]

那么存在一个函数 \(u \in \mathcal{H}\)，使得 \(|u(x)| < |u(a)|\)（对一切 \(x \neq a\)）。

设 \(\left(\mathrm{V}_{n}\right)\) ( \(n \geqslant 1\) ) 是 \(a\) 的一个基本邻域系，并设 \(\lambda, \mu, \varepsilon\) 是满足下列关系的数：

\[
0 <   \lambda <   1, \quad 1 <   \mu <   \mu + \varepsilon \leqslant 1 + \lambda .
\]

于是 \(0 < \lambda / \mu < 1 / \mu < 1\)。我们将对 \(n\) ( \(n \geqslant 1\) ) 用归纳法定义开邻域的一个递减序列 \((\mathrm{U}_n)\)（\(a\) 的），使得 \(\mathrm{U}_n \subset \mathrm{V}_n\)（对一切 \(n\)），以及函数的一个序列 \((h_n)\)（\(\mathcal{H}\) 中的，满足下列关系）：

\[
\left| h _ {n} (x) \right| \leqslant \mu \quad \text { for   all } x \in X\tag{\ensuremath{13_n}}
\]

\[
h _ {n} (a) = 1\tag{\ensuremath{14_n}}
\]

\[
\left| h _ {n} (x) \right| \leqslant \lambda \quad \text { for   all } x \in X - U _ {n}\tag{\ensuremath{15_n}}
\]

\[
\left| \sum_ {j = 1} ^ {n} \lambda^ {j} h _ {j} (y) \right| <   \sum_ {j = 1} ^ {n + 1} \lambda^ {j} \quad \text { for   all } y \in X.\tag{\ensuremath{16_n}}
\]

假定 \(h_m\) 与 \(U_m\)（对 \(1 \leqslant m < n\)）已定义，它们满足前面四个条件（把 \(n\) 换成 \(m\)）；另一方面，令 \(U_0 = X\)。函数 \(\sum_{j=1}^{n-1} \lambda^j h_j\)（当 \(n = 1\) 时等于 0）是连续的，且取值 \(\sum_{j=1}^{n-1} \lambda^j\)（在点 \(a\) 处）；因此存在开邻域 \(U_n\)（\(a\) 的，含于 \(U_{n-1} \cap V_n\)），使得

\[
\left| \sum_ {j = 1} ^ {n - 1} \lambda^ {j} h _ {j} (y) \right| <   \sum_ {j = 1} ^ {n - 1} \lambda^ {j} + \varepsilon \lambda^ {n} \quad \text { for   all } y \in \mathrm{U} _ {n}.\tag{17}
\]

由假设，存在一个函数 \(f \in \mathcal{H}\)，使得

\[
\begin{array}{l} | f (x) | \leqslant 1 \quad \text { for   all } x \in X, \quad | f (a) | \geqslant 1 / \mu , \\ | f (x) | \leqslant \lambda / \mu \quad \text { for } x \in X - U _ {n}. \end{array}
\]

令 \(h_n = f / f(a)\)；于是关系 \((13_n), (14_n)\) 与 \((15_n)\) 成立；令

\[
g = \sum_ {j = 1} ^ {n} \lambda^ {j} h _ {j} = \sum_ {j = 1} ^ {n - 1} \lambda^ {j} h _ {j} + \lambda^ {n} h _ {n}.
\]

由 (17) 与 \((13_{n})\)，对 \(y \in \mathrm{U}_n\) 我们有

\[
| g (y) | <   \sum_ {j = 1} ^ {n - 1} \lambda^ {j} + \varepsilon \lambda^ {n} + \mu \lambda^ {n} \leqslant \sum_ {j = 1} ^ {n + 1} \lambda^ {j},
\]

因为 \(\varepsilon + \mu \leqslant 1 + \lambda\)；对 \(x \in \mathrm{X} - \mathrm{U}_n\)，我们有 \(|h_p(x)| \leqslant \lambda\)（对 \(1 \leqslant p \leqslant n\)），从而也有

\[
| g (x) | \leqslant \sum_ {j = 2} ^ {n + 1} \lambda^ {j} <   \sum_ {j = 1} ^ {n + 1} \lambda^ {j},
\]

这就完成了 \((16_{n})\) 的证明。

这样，级数 \(\sum_{n=1}^{\infty} \lambda^n h_n\) 在 X 中是正规收敛的，因为 \(\lambda < 1\) 且 \(|h_n(x)| \leqslant \mu\)（对一切 \(n\) 及一切 \(x \in \mathbf{X}\)）；用 \(u\) 表示它的和；由于 \(\mathcal{H}\) 是闭的，它属于 \(\mathcal{H}\)。由关系 \((14_n)\)，我们有 \(u(a) = \sum_{n=1}^{\infty} \lambda^n\)；另一方面，若 \(x \neq a\)，则存在一个整数 \(n\) 使得 \(x \notin \mathrm{U}_{n+1}\)；于是 \(|h_{n+k}(x)| \leqslant \lambda\)（对一切 \(k \geqslant 1\)，由关系 \((15_n)\)）；从而，利用 \((16_n)\)，得到

\[
\begin{array}{l} | u (x) | \leqslant \left| \sum_ {j = 1} ^ {n} \lambda^ {j} h _ {j} (x) \right| + \left| \sum_ {j = n + 1} ^ {\infty} \lambda^ {j} h _ {j} (x) \right| <   \sum_ {j = 1} ^ {n + 1} \lambda^ {j} + \lambda \sum_ {j = n + 1} ^ {\infty} \lambda^ {j} \\ = \sum_ {j = 1} ^ {\infty} \lambda^ {j} = | u (a) |. \end{array}
\]

定理 2 (E. Bishop). --- 设 X 是一个紧空间，A 是复 Banach 代数 \(\mathcal{C}(\mathrm{X};\mathbf{C})\) 的一个闭子代数。假定 A 包含常数且分离 X 的点。设 \(a\) 是 X 的一个点；下列条件等价：

\begin{enumerate}
\def\labelenumi{\alph{enumi})}
\item
  存在一个函数 \(f \in \mathcal{A}\)，使得 \(|f(x)| < |f(a)|\)（对一切 \(x \neq a\)）。
\item
  点 \(a\) 是 \(\mathcal{A}\)-极值点且具有可数的基本邻域系。
\end{enumerate}

\(a) \Rightarrow b)\)：设 \(f \in \mathcal{A}\) 满足 \(|f(a)| > |f(x)|\)（对 \(x \neq a\)）；由 No.~4 的命题 9，\(a\) 是 \(\mathcal{A}\)-极值点。另一方面，若 \(\mathrm{U}_n\) 是由 \(x \in \mathrm{X}\)（满足 \(|f(x)| > |f(a)| - 1/n\)）组成之集，则 \(\mathrm{U}_n\) 是 \(a\) 的一个开邻域，且诸 \(\mathrm{U}_n\) 的交化为 \(a\)；由于 \(\mathrm{X}\) 是紧的，诸 \(\mathrm{U}_n\) 构成 \(a\) 的一个基本邻域系（GT, I, §9, No.~1, 定理 1）。

\(b) \Rightarrow a)\)：只须验证 \(b)\) 蕴涵引理 4 的假设。采用该引理的记号，令 \(\varepsilon = \log d / \log c\)；于是 \(0 < \varepsilon < 1\)。由于 \(a\) 是 \(\mathcal{A}_r\)-极值点，存在一个函数 \(g \in \dot{\mathcal{A}}\)，使得

\[
\mathcal {R} (g) \geqslant 0, \quad \mathcal {R} (g (a)) \leqslant \varepsilon , \quad \mathcal {R} (g (x)) \geqslant 1 \text {for} x \in X - U
\]

（No.~3, 命题 6, b)）。令 \(f = c^g\)；由于 \(f\) 是正规收敛级数 \(\sum_{n=0}^{\infty} (\log c)^n g^n / n!\) 的和，我们有 \(f \in \mathcal{A}\)，且

\[
| f | \leqslant 1, \quad | f (a) | \geqslant c ^ {\varepsilon} = d, \quad | f (x) | \leqslant c \text { for } x \in X - U.\tag{Q.E.D.}
\]

推论. --- 再假定 X 是可度量化的。那么下列性质等价：

\begin{enumerate}
\def\labelenumi{\alph{enumi})}
\item
  \(a\) 是 X 的一个 \(\mathcal{A}\)-极值点。
\item
  存在 \(u \in \mathcal{A}\)，使得 \(|u(x)| < |u(a)|\)（对一切 \(x \neq a\)）。
\item
  设 \(\mathfrak{M}\) 是子集 \(M\)（\(X\) 的）所成之集：对每个函数 \(f \in \mathcal{A}\)，\(|f|\) 在 \(X\) 中的上确界至少在 \(M\) 的一点处达到。那么 a 属于所有的集 \(M \in \mathfrak{M}\)。
\item
  设 \(\mathfrak{N}\) 是 X 的这样的子集 N 所成之集：对每个函数 \(f \in \mathcal{A}\)，\(\mathcal{R}(f)\) 在 X 中的上确界至少在 N 的一点处达到。那么 a 属于所有的集 \(N \in \mathfrak{N}\)。
\end{enumerate}

换言之，

\[
\operatorname{Ch} _ {\mathcal {A}} (\mathrm{X}) = \bigcap_ {\mathrm{M} \in \mathfrak {M}} \mathrm{M} = \bigcap_ {\mathrm{N} \in \mathfrak {N}} \mathrm{N}.\tag{18}
\]

由于在可度量化空间中每个点都具有可数的基本邻域系，\(a\) 与 \(b\) 的等价性由定理 2 得出。我们来证明 \(b\) 蕴涵 \(c\)：事实上，\(a\) 是 \(|u|\) 达到其上确界的唯一点；另一方面，\(c\) 蕴涵 \(a\)，因为对每个 \(f \in \mathcal{A}\)，\(\mathrm{Ch}_{\mathcal{A}}(\mathrm{X})\) 与使 \(|f|\) 达到其上确界的点所成之集相交（No.~4, 命题 9）。同样的推理，利用 No.~3 的命题 5，表明 \(d\) 蕴涵 \(a\)）。最后，为看到 \(b\) 蕴涵 \(d\)，我们可以限于 X 不化为单独一点 \(a\) 的情形，因此 \(u(a) \neq 0\)；于是函数 \(v = u / u(a)\) 属于 \(\mathcal{A}\)，且我们有 \(v(a) = 1\) 以及 \(|v(x)| < 1\)（对 \(x \neq a\)），从而 \(\mathcal{R}(v(a)) = 1\) 且 \(\mathcal{R}(v(x)) < 1\)（对 \(x \neq a\)）。由于函数 \(\mathcal{R}(v)\) 仅在点 \(a\) 处达到其上确界，的确有 \(a \in N\)（对每个 \(N \in \mathfrak{N}\)）。

*例. --- 设 \(\mathrm{X}_1\) 是由点 \((z_1, z_2) \in \mathbf{C}^2\)（满足 \(|z_1|^2 + |z_2|^2 \leqslant 1\)）组成之集（\(\mathbf{R}^4\) 中的单位球），并设 \(\mathcal{A}_1'\) 是全纯函数在 \(\mathrm{X}_1\) 上的限制所成之集，这些全纯函数取值于 \(\mathbf{C}\)、定义在 \(\mathrm{X}_1\) 于 \(\mathbf{C}^2\) 中的一个邻域内（所论邻域依所考虑的函数而定）；设 \(\mathcal{A}_1\) 是 \(\mathcal{A}_1'\) 在 \(\mathcal{C}(\mathrm{X}_1; \mathbf{C})\) 中的闭包，它显然是 \(\mathcal{C}(\mathrm{X}_1; \mathbf{C})\) 的一个闭复子代数且分离 \(\mathrm{X}_1\) 的点。对全纯函数应用“极大值原理”表明 \(\mathrm{Ch}_{\mathcal{A}_1}(\mathrm{X}_1)\) 是球面 \(\mathbf{S}_3\)。

在前面的定义中，把 \(X_1\) 换成“多圆柱”\(X_2\)（由关系 \(|z_1| \leqslant 1\) 与 \(|z_2| \leqslant 1\) 定义），这就给出子代数 \(\mathscr{A}_2'\) 与 \(\mathscr{A}_2\)（\(\mathcal{A}_2^{\prime}\) 的闭包，属于 \(\mathcal{C}(\mathrm{X}_2; \mathbf{C})\)）。这里，极大值原理表明 \(\mathrm{Ch}_{\mathcal{A}_2}(\mathrm{X}_2)\) 是由关系 \(|z_1| = 1\) 与 \(|z_2| = 1\) 定义的“环面”。

由这些结果可以推出：不存在 \(X_{1}\) 的一个开邻域到 \(X_{2}\) 的一个开邻域上、且把 \(X_{1}\) 变成 \(X_{2}\) 的解析同构；因为，若 v 是这样一个映射在 \(X_{1}\) 上的限制，则将有 \(A_{2}=vA_{1}v^{-1}\)，于是 v 将把 \(S_{3}\) 变成一个同胚于 \(T^{2}\) 的空间，这是荒谬的，因为 \(S_{3}\) 是单连通的而 \(T^{2}\) 不是。然而人们会注意到，空间 \(X_{1}\) 与 \(X_{2}\) 是同胚的，二者都是 \(R^{4}\) 中具有非空内部的有界凸集。

\subsection{积分表示的唯一性}

设 E 是一个 Hausdorff 弱局部凸空间（TVS, II, §6, No.~2），C 是 E 中一个真尖凸锥。已知 C 是 E 中某个与 E 的向量空间结构相容的序关系下的 \(\geqslant0\) 元素组成之集。当说 C 是格序的时候，当然理解为 E 的序在 C 上诱导的序。

引理 5. --- 假定 C 是弱完备的。设 A 是 E 上连续线性形式在 C 上的限制所成之集。设 \((f_{\lambda})_{\lambda \in \Lambda}\) 是 A 中元素的一个有限族，且 \(f = \sup(f_{\lambda})\)。对每个 \(x \in C\)，定义

\[
\overline {{{f}}} (x) = \sup \left(f (x _ {1}) + f (x _ {2}) + \dots + f (x _ {n})\right),
\]

其中上确界取遍集 \(S_{x}\)——由 C 中元素的序列 \((x_{1}, x_{2}, \ldots, x_{n})\)（满足 \(x_{1} + x_{2} + \cdots + x_{n} = x\)）组成。设 \(\operatorname{Card}\Lambda = p\)。那么存在 \((y_{1}, \ldots, y_{p}) \in \operatorname{S}_{x}\) 使得 \(\overline{f}(x) = f(y_{1}) + \cdots + f(y_{p})\)。

用 \(f_1, f_2, \ldots, f_p\) 表示族 \((f_\lambda)\) 的元素。对 \(k = 1, 2, \ldots, p\)，设 \(C_k\) 是由满足下列关系的 \(y \in C\) 组成之集：

\[
f _ {1} (y) <   f (y), f _ {2} (y) <   f (y), \dots , f _ {k - 1} (y) <   f (y), f _ {k} (y) = f (y).
\]

诸 \(C_{k}\) 是两两不相交的凸锥，其并为 C。设 C 中的 \(x_{1}, x_{2}, \ldots, x_{n}\) 满足 \(x_{1} + x_{2} + \cdots + x_{n} = x\)。设 \(y_{k}\) 是那些 \(x_{i}\)（属于 \(C_{k}\) 的）之和。那么 \(y_{1} + y_{2} + \cdots + y_{p} = x\)。由于 f 在 \(C_{k}\) 上是仿射的，故 \(f(y_{1}) + \cdots + f(y_{p}) = f(x_{1}) + \cdots + f(x_{n})\)。因此

\[
f (x) = \sup \bigl (f (y _ {1}) + \dots + f (y _ {p}) \bigr),\tag{19}
\]

其中 \((y_{1}, y_{2}, \ldots, y_{p})\) 取遍由满足 \(y_{1} + y_{2} + \cdots + y_{p} = x\) 的 C 中 p 个点的序列组成之集。令 \(D = C \cap (x - C)\)。由于 D 是紧的（TVS, II, §6, No.~8, 命题 11 的推论 2），元素 \((y_{1}, \ldots, y_{p})\)（\(D^{p}\) 中的，满足 \(y_{1} + \cdots + y_{p} = x\)）所成之集也是紧的，因而上确界 (19) 可以达到。

引理 6. --- 我们保持引理 5 的假设与记号，并假定诸 \(f_{\lambda}\) 是正的。函数 \(\overline{f}\) 在 C 中是正齐次的、凹的且上半连续的。若 C 是格序的，则它是仿射的。

显然 \(\overline{f}\) 是正齐次的。设 x, y 属于 C。若 C 中的 \(x_{1}, \ldots, x_{m}, y_{1}, \ldots, y_{n}\) 满足 \(x_{1} + \cdots + x_{m} = x\)，\(y_{1} + \cdots + y_{n} = y\)，则 \(x_{1} + \cdots + x_{m} + y_{1} + \cdots + y_{n} = x + y\)，因此

\[
f (x _ {1}) + \dots + f (x _ {m}) + f (y _ {1}) + \dots + f (y _ {n}) \leqslant \overline {{f}} (x + y);
\]

由此得出 \(\overline{f}(x) + \overline{f}(y) \leqslant \overline{f}(x + y)\)，因此 \(\overline{f}\) 是凹的。设 \(L\)（相应地，\(L_{\lambda}\)）是由 \((t, x) \in \mathbf{R} \times E\)（满足 \(x \in C\) 且 \(0 \leqslant t \leqslant \overline{f}(x)\)（相应地，\(0 \leqslant t \leqslant f_{\lambda}(x)\)））组成之集。每个 \(L_{\lambda}\) 都在弱完备真凸锥 \(\mathbf{R}_{+} \times C\) 中是闭的，因此和 \(\sum_{\lambda \in \Lambda} L_{\lambda}\) 是闭的（TVS, II, §6, No.~8, 命题 11 的推论 2）。由引理 5，这个和等于 L。于是 L 是闭的，这就证明 \(\overline{f}\) 是上半连续的。最后，假定 C 是格序的，我们来证明 \(\overline{f}\) 是凸的。设 \(x, y\) 属于 C 且 \(\varepsilon > 0\)。存在 C 中的 \(z_1, z_2, \ldots, z_n\)，使得 \(f(z_1) + \cdots + f(z_n) \geqslant \overline{f}(x + y) - \varepsilon\) 且 \(z_1 + \cdots + z_n = x + y\)。向量空间 C - C 对 E 的序所诱导的序是格序的（A, VI, §1, No.~9, 命题 8）。由分解定理（同前，No.~10, 定理 1），存在 C 中的 \(x_1, \ldots, x_n, y_1, \ldots, y_n\)，使得

\[
x _ {1} + y _ {1} = z _ {1}, \dots , x _ {n} + y _ {n} = z _ {n}, x _ {1} + \dots + x _ {n} = x, y _ {1} + \dots + y _ {n} = y.
\]

于是，由于 \(f\) 是正齐次且凸的，

\[
\begin{array}{r l} \overline {{f}} (x + y) & \leqslant \varepsilon + f (z _ {1}) + \dots + f (z _ {n}) \\ & \leqslant \varepsilon + f (x _ {1}) + f (y _ {1}) + \dots + f (x _ {n}) + f (y _ {n}) \\ & \leqslant \varepsilon + \overline {{f}} (x) + \overline {{f}} (y). \end{array}
\]

由于 \(\varepsilon\) 是任意的数 \(>0\)，我们已证明 \(\overline{f}\) 的确是凸的。

定理 3 (Choquet). --- 设 E 是一个 Hausdorff 弱局部凸空间，C 是 E 中以 0 为顶点的弱完备真凸锥，G 是 C 的极值母线之并，K 是 C 的一个紧凸子集，\(\lambda\) 与 \(\lambda'\) 是 K 上质量为 1、具有相同重心的正测度，且满足 \(\lambda^{*}(\mathrm{K} - (\mathrm{K} \cap \mathrm{G})) = \lambda'^{*}(\mathrm{K} - (\mathrm{K} \cap \mathrm{G})) = 0\)。假定 C 是格序的。那么，对 C 上每个下半连续、正齐次的凸函数 \(f \geqslant 0\)，有 \(\lambda^{*}(f|K) = \lambda'^{*}(f|K)\)。

设 \(\mathcal{A}\)（相应地，\(\mathcal{A}'\)）是 E 上连续线性形式（相应地，仿射函数）在 \(\dot{\mathrm{C}}\) 上的限制所成之集。我们知道（TVS, II, §5, No.~4, 注 2）\(f\) 是 \(\mathcal{A}\) 中 \(\leqslant f\) 的元素所成之集的上包络。形如 \(\sup(f_1, \ldots, f_p)\) 的函数所成之集（其中 \(f_1, \ldots, f_p\) 属于 \(\mathscr{A}\)，\(f_1 \geqslant 0, \ldots, f_p \geqslant 0\)）是一个递增有向集，且以 \(f\) 为其上包络。考虑到 §1, No.~1, 定理 1，只须验证等式 \(\lambda(f|K) = \lambda'(f|K)\)（当 \(f\) 具有前述形式时）。

按引理 5 那样定义 \(\overline{f}\)。显然，\(\overline{f}(y) = f(y)\)（若 \(y \in G\)）。由于 \(\lambda^{*}(\mathrm{K} - (\mathrm{K} \cap \mathrm{G})) = 0\)，我们有 \(\lambda(f|K) = \lambda(\overline{f}|K)\)。由引理 6，\(\overline{f}\) 是仿射的且上半连续。因此 \(\overline{f}|K\) 是 \(\mathcal{A}'\) 的元素在 K 上的限制所成的一个递减有向集的下包络（TVS, II, §5, No.~4, 命题 6）。设 \(x \in K\) 是 \(\lambda\) 的重心。若 \(g \in \mathcal{A}\)，则 \(\lambda(g|K) = g(x)\)。因此 \(\lambda(\overline{f}|K) = \overline{f}(x)\)（§4, No.~4, 命题 5 的推论 2）。于是 \(\lambda(f|K) = \overline{f}(x)\)，并且同理可见 \(\lambda'(f|K) = \overline{f}(x)\)。

推论. --- 设 E 是一个 Hausdorff 局部凸空间，C 是 E 中以 0 为顶点、具有紧底 M 的真凸锥，并设 G 是 C 的极值母线之并。设 \(x \in M\)。若 C 是格序的，则 M 上至多存在一个质量为 1 的正测度 \(\lambda\)，使得 \(\lambda^{*}(\mathrm{M} - (\mathrm{G} \cap \mathrm{M})) = 0\)，且以 x 为重心。

把 E 的拓扑换成弱化的拓扑（这不改变 M 的拓扑），就可以假定 E 是弱空间。设 \(\lambda\) 与 \(\lambda'\) 是 M 上具有所述性质的两个测度，并设 h 是 E 上使 M 为 C 与方程 \(h(x)=1\) 的超平面之交的连续线性形式。设 S 是 \(\mathcal{C}(\mathrm{M})\) 中由 C 上正齐次且连续的 \(\geqslant0\) 凸函数在 M 上的限制组成的子集。锥 C 是弱完备的（TVS, II, §7, No.~3）。由定理 3，\(\lambda(f)=\lambda'(f)\)（对每个 \(f\in S\)）。

若 \(f_{1}, f_{2}, f_{3}, f_{4}\) 属于 \(\mathcal{S}\)，则

\[
\sup \left(f _ {1} - f _ {2}, f _ {3} - f _ {4}\right) = \sup \left(f _ {1} + f _ {4}, f _ {3} + f _ {2}\right) - \left(f _ {2} + f _ {4}\right) \in \mathscr {S} - \mathscr {S}
\]

\[
\inf \left(f _ {1} - f _ {2}, f _ {3} - f _ {4}\right) = - \sup \left(f _ {2} - f _ {1}, f _ {4} - f _ {3}\right) \in \mathscr {S} - \mathscr {S}.
\]

由于 \(h|M \in \mathcal{S}\)，\(\mathcal{S} - \mathcal{S}\) 包含常值函数。若 \(x\) 与 \(y\) 是 \(M\) 的两个不同点，则存在 \(E\) 上的一个连续线性形式，它分离 \(x\) 与 \(y\)，且这一形式是两个在 \(C\) 上为正的连续线性形式之差（TVS, II, §6, No.~8, 引理 1）。由上述事实推出：对实数 \(\alpha, \beta\)，存在 \(f \in \mathcal{S} - \mathcal{S}\) 使得 \(f(x) = \alpha\)，\(f(y) = \beta\)。于是 \(\mathcal{S} - \mathcal{S}\) 在 \(\mathcal{C}(M)\) 中对一致收敛拓扑是稠密的（GT, X, §4, No.~1, 命题 2 的推论）。由于 \(\lambda\) 与 \(\lambda'\) 在 \(\mathcal{S} - \mathcal{S}\) 上一致，我们有 \(\lambda = \lambda'\)。

\chapter*{习题}
\phantomsection
\addcontentsline{toc}{chapter}{习题}
\setcounter{exercisegroup}{1}
\edef\CurrentExerciseGroup{\arabic{exercisegroup}}
\section*{\texorpdfstring{\S\CurrentExerciseGroup}{第{\CurrentExerciseGroup}节}}
\phantomsection
\addcontentsline{toc}{section}{第{\CurrentExerciseGroup}节习题}

\begin{enumerate}
\def\labelenumi{\arabic{enumi})}
\tightlist
\item
  试证：若 \(f\) 与 \(g\) 是两个数值函数 \(\geqslant 0\)（定义在 \(X\) 上），而 \(\mu\) 是 \(X\) 上的一个正测度，则
\end{enumerate}

\[
\mu^ {*} (\sup (f, g)) + \mu^ {*} (\inf (f, g)) \leqslant \mu^ {*} (f) + \mu^ {*} (g).
\]

由此推出：若 A 与 B 是 X 的任意子集，则

\[
\mu^ {*} (\mathrm{A} \cup \mathrm{B}) + \mu^ {*} (\mathrm{A} \cap \mathrm{B}) \leqslant \mu^ {*} (\mathrm{A}) + \mu^ {*} (\mathrm{B})
\]

（参见 §4, 习题 8 d)）。

\begin{enumerate}
\def\labelenumi{\arabic{enumi})}
\setcounter{enumi}{1}
\item
  对每个 \(x \in \mathbf{X}\)，试证对每个数值函数 \(f \geqslant 0\)（定义在 \(\mathbf{X}\) 上），有 \(\varepsilon_x^*(f) = f(x)\)。
\item
  在 \(\mathbf{R}\) 中给出一个不是相对紧、但其外测度（对 Lebesgue 测度而言）有限的开集的例子。
\item
  设 \(\mathrm{X}\) 是一个局部紧空间，\(\alpha\) 是一个有限数值函数 \(\geqslant 0\)（在 \(\mathrm{X}\) 上），使得对每个紧子集 \(\mathrm{K}\)（\(\mathrm{X}\) 的），\(\sum_{x\in \mathrm{K}}\alpha (x)\) 是有限的；设 \(\mu\) 是正测度
\end{enumerate}

它由质量 \(\alpha(x)\) 定义在 X 上（Ch. III, §1, No.~3, 例 I）。

\begin{enumerate}
\def\labelenumi{\alph{enumi})}
\tightlist
\item
  试证：对每个函数 \(f \geqslant 0\)（在 \(X\) 上下半连续），有 \(\mu^{*}(f) = \sum_{x \in X} \alpha(x)f(x)\)，这里约定 \(\alpha(x)f(x) = 0\) 当 \(\alpha(x) = 0\) 且 \(f(x) = -\infty\)，或者当
\end{enumerate}

\[
f (x) = + \infty
\]

\begin{enumerate}
\def\labelenumi{\alph{enumi})}
\setcounter{enumi}{1}
\tightlist
\item
  时同样成立。设 \(f \geqslant 0\) 是任一数值函数（定义在 \(X\) 上）。试证：若 \(\mu^{*}(f) < +\infty\)，则 \(\mu^{*}(f) = \sum_{x \in X} \alpha(x)f(x)\)，其中约定与 a) 中相同。（参见习题 5。）
\end{enumerate}

¶ 5) 设 X 是平面 \(\mathbf{R}^2\) 的由直线 \(D = \{0\} \times \mathbf{R}\) 与点 \((1/n, k/n^2)\) 所成之集的并组成的子集，其中 \(n\) 取遍整数 \(>0\) 的集，\(k\) 取遍有理整数集 \(\mathbf{Z}\)。

\begin{enumerate}
\def\labelenumi{\alph{enumi})}
\item
  对 D 的每个点 \((0,y)\) 及每个整数 n \textgreater{} 0，设 \(\mathrm{T}_{n}(y)\) 是 X 中点 \((u,v)\)（满足 \(u \leqslant 1/n\) 且 \(|v - y| \leqslant u\)）组成之集。试证：若对 D 中每个点 \((0,y)\) 取诸 \(\mathrm{T}_{n}(y)\) (n \textgreater{} 0) 之集作为基本邻域系，并取化为该点本身的唯一集作为 X 的其余每个点的基本邻域系，就在 X 上定义了一个拓扑 T，对它而言 X 是一个局部紧但非仿紧的空间。
\item
  设 \(\alpha\) 是 \(X\) 上的数值函数，它在 \(D\) 上等于 0，且等于 \(1/n^3\)（在每个点 \((1/n, k/n^2)\) 处）。试证：对每个紧子集 \(K\)（\(X\) 的），\(\sum_{x \in K} \alpha(x)\)
\end{enumerate}

有限；设 \(\mu\) 是 X 上由质量 \(\alpha(x)\) 定义的正测度。试证 \(\mu^{*}(D) = +\infty\)，尽管在 D 上 \(\alpha(x) = 0\)。（若关于 \(\mathcal{T}\) 的一个开集 U 包含 D，试证：在 D 上存在一个不化为一点的区间 \(a \leqslant y \leqslant b\)、一个在该区间中稠密的集 B（关于 \(\mathbf{R}\) 的通常拓扑）及一个整数 \(n > 0\)，使得对每个 \(y \in B\)，有 \(\mathrm{T}_{n}(y) \subset \mathrm{U}\)；为此可利用 Baire 定理。）

\begin{enumerate}
\def\labelenumi{\arabic{enumi})}
\setcounter{enumi}{5}
\tightlist
\item
  设 \(f\) 是一个数值函数 \(\geqslant 0\)（定义在 \(X\) 上）。
\end{enumerate}

\begin{enumerate}
\def\labelenumi{\alph{enumi})}
\item
  试证：要使 \(\mu \mapsto \mu^{*}(f)\) 从 \(\mathcal{M}_{+}(X)\) 到 \(\overline{\mathbf{R}}\) 的映射关于淡拓扑连续，必须（且只须）\(f\) 是具有紧支集的连续函数（利用习题 2）。要使 \(\mu \mapsto \mu^{*}(f)\) 关于淡拓扑下半连续，必须（且只须，参见命题 4）\(f\) 下半连续。
\item
  试证：要使 \(\mu \mapsto \mu^{*}(f)\) 从 \(\mathcal{M}_{+}(\mathrm{X})\) 到 \(\overline{\mathbf{R}}\) 的映射关于拟强拓扑（Ch. III, §1, 习题 8）连续，必须且只须 \(f\) 有界且具有紧支集（方法与 \(a\) ) 的类似）。由此推出，对每个在可数个紧集之并的余集上为零的函数 \(f \geqslant 0\)，映射 \(\mu \mapsto \mu^{*}(f)\) 关于拟强拓扑下半连续（利用定理 3）（参见习题 7 b)）。
\item
  试证：要使 \(\mu \mapsto \mu^{*}(f)\) 从 \(\mathcal{M}_{+}(\mathrm{X}) \cap \mathcal{M}^{1}(\mathrm{X})\) 到 \(\overline{\mathbf{R}}\) 的映射关于超强拓扑（Ch. III, §1, 习题 15）连续，必须且只须 \(f\) 有界；对每个函数 \(f \geqslant 0\)（定义在 \(\mathrm{X}\) 上），\(\mu \mapsto \mu^{*}(f)\) 关于超强拓扑下半连续。
\item
  试证：要使 \(\mu \mapsto \mu^{*}(f)\) 从 \(\mathcal{M}_{+}(\mathrm{X}) \cap \mathcal{M}^{1}(\mathrm{X})\) 到 \(\overline{\mathbf{R}}\) 的映射关于弱拓扑（Ch. III, §1, 习题 15）连续，必须且只须 \(f\) 在 \(\mathrm{X}\) 上连续且在无穷远点趋于 0；要使 \(\mu \mapsto \mu^{*}(f)\) 关于弱拓扑下半连续，必须且只须 \(f\) 下半连续。
\end{enumerate}

\begin{enumerate}
\def\labelenumi{\arabic{enumi})}
\setcounter{enumi}{6}
\tightlist
\item
  \begin{enumerate}
  \def\labelenumii{\alph{enumii})}
  \tightlist
  \item
    设 \((\mu_{n})\) 是局部紧空间 X 上一个递增的测度序列 \(\geqslant0\)；假设该序列在 \(\mathcal{M}_{+}(X)\) 中有上界，并用 \(\mu\) 表示它的上确界。设 f 是定义在 X 上且在可数个紧集之并的余集上为零的函数 \(\geqslant0\)。试证 \(\mu^{*}(f)=\sup\mu_{n}^{*}(f)\)（参见习题 6 b)）。
  \end{enumerate}
\end{enumerate}

\begin{enumerate}
\def\labelenumi{\alph{enumi})}
\setcounter{enumi}{1}
\tightlist
\item
  设 X 与 \(\mu\) 是习题 5 中定义的局部紧空间与测度。设 \(\alpha_{n}\) 是这样的数值函数：它由 \(\alpha\)（习题 5 的记号）在每个点 \((1/m, k/m^{2})\)（这里 \(m \leqslant n\)）处的值给出，而在 X 的其余各点处等于 0；设 \(\mu_{n}\) 是由质量 \(\alpha_{n}(x)\) 定义的测度。试证：\(\mu\) 是序列 \((\mu_{n})\) 在 \(\mathcal{M}_{+}(X)\) 中的上确界，且对一切 n 有 \(\mu_{n}^{*}(D)=0\)，但 \(\mu^{*}(D)=+\infty\)。
\end{enumerate}

\begin{enumerate}
\def\labelenumi{\arabic{enumi})}
\setcounter{enumi}{7}
\tightlist
\item
  \begin{enumerate}
  \def\labelenumii{\alph{enumii})}
  \tightlist
  \item
    设 \((\mu_n)\) 是一个递减的测度序列 \(\geqslant 0\)（在局部紧空间 \(X\) 上），并设 \(\mu\) 是该序列在 \(\mathcal{M}_+(X)\) 中的下确界。试证：若 \(f\) 是一个函数 \(\geqslant 0\)，使得从某个指标起 \(\mu_n^*(f) < +\infty\)，则 \(\mu^*(f) = \inf_n \mu_n^*(f)\)（当 \(g\) 下半连续、为正且满足 \(\mu^*(g) < +\infty\) 时，注意存在一个序列 \((h_m)\)，它由具有紧支集的连续函数 \(\geqslant 0\) 组成，使得 \(\sum_{m=1}^{\infty} h_m \leqslant g\) 且 \(\mu_n^*(g) = \sum_{m=1}^{\infty} \mu_n^*(h_m)\) 对每个指标 \(n\)。
  \end{enumerate}
\end{enumerate}

\begin{enumerate}
\def\labelenumi{\alph{enumi})}
\setcounter{enumi}{1}
\tightlist
\item
  在离散空间 X = N 上，设 \(\mu_{n}\) 是在每个点 \(m \geqslant n\) 处放置质量 +1 而定义的测度。试证：递减序列 \((\mu_{n})\) 在 \(\mathcal{M}_{+}(X)\) 中的下确界为 0，但对一切 n 有 \(\mu_{n}^{*}(X) = +\infty\)。
\end{enumerate}

\setcounter{exercisegroup}{3}
\edef\CurrentExerciseGroup{\arabic{exercisegroup}}
\section*{\texorpdfstring{\S\CurrentExerciseGroup}{第{\CurrentExerciseGroup}节}}
\phantomsection
\addcontentsline{toc}{section}{第{\CurrentExerciseGroup}节习题}

\begin{enumerate}
\def\labelenumi{\arabic{enumi})}
\item
  设 \(\mu\) 是区间 \(X = [0,1[\)（\(\mathbf{R}\) 内）上的 Lebesgue 测度。对每个整数 \(n = 2^h + k\)（\(h \geqslant 0\)，\(0 \leqslant k < 2^h\)），设 \(f_n\) 是在区间 \([k/2^h, (k+1)/2^h[\) 上等于 1、在 \(X\) 中其余各处等于 0 的函数。试证：序列 \((f_n)\) 以 \(p\) 阶平均收敛于 0（对每个 \(p > 0\)），但序列 \((f_n(x))\) 对任何点 \(x \in X\) 都不收敛。
\item
  试证：每个数值函数 \(f\)（属于 \(\mathcal{L}^p(\mathrm{X};\mathbf{R})\)）都几乎处处等于差 \(g_1 - g_2\)（即属于 \(\mathcal{L}^p(\mathrm{X};\mathbf{R})\) 的两个下半连续正函数之差）（注意 \(f(x)\) 几乎处处等于一个绝对收敛级数 \(\sum_{n=1}^{\infty} f_n(x)\) 的和，其中诸 \(f_n\) 是具有紧支集的连续函数）。
\item
  设 \(\mathbf{X}\) 是一个局部紧空间，\(\alpha\) 是一个有限数值函数 \(\geqslant 0\)，定义在 \(\mathbf{X}\) 上，使得对每个紧子集 \(\mathbf{K}\)（含于 \(\mathbf{X}\)）有 \(\sum_{x\in \mathbf{K}}\alpha (x) < +\infty\)。设 \(\mu\) 是 \(\mathbf{X}\) 上由质量 \(\alpha (x)\) 定义的测度。试证：对这个测度，\(\mathcal{F}_{\mathrm{F}}^{p} = \mathcal{L}_{\mathrm{F}}^{p}\) 对每个有限的 \(p\geqslant 1\) 与每个 Banach 空间 \(\mathbf{F}\) 成立。
\end{enumerate}

\setcounter{exercisegroup}{4}
\edef\CurrentExerciseGroup{\arabic{exercisegroup}}
\section*{\texorpdfstring{\S\CurrentExerciseGroup}{第{\CurrentExerciseGroup}节}}
\phantomsection
\addcontentsline{toc}{section}{第{\CurrentExerciseGroup}节习题}

\begin{enumerate}
\def\labelenumi{\arabic{enumi})}
\tightlist
\item
  设 H 是可积函数所成的一个集合，按关系 \(\leqslant\) 有向，这些函数 \(\geqslant 0\) 且满足 \(\sup N_{1}(f) < +\infty\)。设 g 是 H 的上包络；为使 g 可积且 \(f \in H\)
\end{enumerate}

在 \(L^{1}\) 中，g 的类 \(\widetilde{g}\) 是诸函数 \(f \in H\) 的类所成有向集的截口滤子之极限，这必须且只须：对每个 \(\varepsilon > 0\)，存在一个函数 \(f_{1} \in H\)、一个由下半连续可积函数组成且按 \(\leqslant\) 有向的集 B、以及一个映射 \(f \mapsto f^{*}\)，它定义在集 \(H_{1}\)（由诸函数 \(f \in H\) 中 \(\leqslant f_{1}\) 者组成）上且取值于集 B，使得 \(f \leqslant f^{*}\) 且 \(\mathrm{N}_{1}(f^{*} - f) \leqslant \varepsilon\) 对每个函数 \(f \in H_{1}\) 成立（利用 No.~4 的定理 3）。

\begin{enumerate}
\def\labelenumi{\arabic{enumi})}
\setcounter{enumi}{1}
\item
  设 \(\mu\) 是 \(\mathbf{R}\) 上的 Lebesgue 测度，\(\Omega\) 是在空间 \(\mathcal{L}^1\) 上赋予 \(\mathbf{R}\) 中逐点收敛拓扑所得的拓扑空间。试证：映射 \(f\mapsto \int fd\mu\)（从 \(\Omega\) 到 \(\mathbf{R}\)）在 \(\Omega\) 的任何点处都不连续。
\item
  设 I 是 \(\mathbf{R}\) 中的一个区间，\(\mathbf{f}\) 是 \(X \times I\) 到一个 Banach 空间 \(F\) 中的映射，满足：\(1^{\circ}\) 对每个 \(\alpha \in I\)，映射 \(t \mapsto \mathbf{f}(t, \alpha)\)（从 \(X\) 到 \(F\)）可积；\(2^{\circ}\) 对每个 \(t \in X\)，映射 \(\alpha \mapsto \mathbf{f}(t, \alpha)\) 有导数 \(\mathbf{f}_{\alpha}'(t, \alpha)\)（在 \(I\) 中）；\(3^{\circ}\) 存在一个可积函数 \(g \geqslant 0\)，使得 \(|\mathbf{f}_{\alpha}'(t, \alpha)| \leqslant g(t)\)（对一切 \(t \in X\) 与一切 \(\alpha \in I\)）。在这些条件下，试证：函数 \(\mathbf{u}(\alpha) = \int \mathbf{f}(t, \alpha) d\mu(t)\) 在 \(I\) 中可微，并且
\end{enumerate}

\[
\mathbf {u} ^ {\prime} (\alpha) = \int \mathbf {f} _ {\alpha} ^ {\prime} (t, \alpha) d \mu (t).
\]

\begin{enumerate}
\def\labelenumi{\arabic{enumi})}
\setcounter{enumi}{3}
\tightlist
\item
  设 \(\mu\) 是区间 \(X = [0,1]\) 上的 Lebesgue 测度。
\end{enumerate}

\begin{enumerate}
\def\labelenumi{\alph{enumi})}
\item
  在 X 中定义一个无处稠密的完全集 A，其测度为任意数 \(\alpha\)，满足 \(0 \leqslant \alpha < 1\)（利用 Cantor 三分集的构造方法）。
\item
  在 X 中定义一个由两两不交的无处稠密可积集组成的序列 \((\mathbf{A}_n)\)，使得 \(\mu (\mathrm{A}_n) = 2^{-n}\)，且使得 \(\mathbf{B}_n = \bigcup_{k = 1}^{n}\mathbf{A}_k\) 的每个余区间都含有 \(\mathrm{A}_{n + 1}\) 的一个测度 \(>0\) 的子集。若 \(\mathbf{A} = \bigcup_{n}\mathbf{A}_n\)，试证：\(\mathbf{A}\) 是测度为 1 的贫集，而 \(\mathbb{C}\mathbf{A}\) 是可忽略的非贫集。
\item
  设 \(\mathrm{H} = \bigcup_{n=0}^{\infty} \mathrm{A}_{2n+1}\)；试证：对每个开区间 \(\mathrm{I} \subset \mathrm{X}\)，I 与 H 以及与 CH 的交都有测度 \textgreater{} 0。设 f 是几乎处处等于特征函数 \(\varphi_{H}\) 的函数；试证：不存在任何由 X 上连续函数组成、在 X 的每点收敛且极限等于 f 的序列 \((f_{n})\)（注意 f 必然在 X 的每点不连续，并利用 GT, IX, §5 的习题 22）。
\end{enumerate}

\begin{enumerate}
\def\labelenumi{\arabic{enumi})}
\setcounter{enumi}{4}
\tightlist
\item
  设 \(\mu\) 是局部紧空间 \(X\) 上的一个正测度。对每个数值函数 \(f\)（无论有限与否、符号如何），只要它定义在 \(X\) 上，我们就用 \(\mu^{*}(f)\)（并称之为 \(f\) 的上积分）表示诸数 \(\mu(h)\)（对满足 \(h \geqslant f\) 的可积下半连续函数 h 而言）的下确界，当这样的函数存在时；在相反情形用 \(+\infty\)；当 \(f \geqslant 0\) 时，这个定义与 §1, No.~3 中的定义一致。我们用 \(\mu_{*}(f)\)（并称之为 \(f\) 的下积分）表示数 \(-\mu^{*}(-f)\)。
\end{enumerate}

\begin{enumerate}
\def\labelenumi{\alph{enumi})}
\item
  试证：若 \(f_{1}\) 与 \(f_{2}\) 是两个数值函数，使得几乎处处 \(f_{1}(x) \leqslant f_{2}(x)\)，则 \(\mu^{*}(f_{1}) \leqslant \mu^{*}(f_{2})\)。
\item
  设 \(f_{1}\) 与 \(f_{2}\) 是两个数值函数，使得 \(\mu^{*}(f_{1}) + \mu^{*}(f_{2})\) 有定义且 \(< +\infty\)；试证：\(f_{1}(x) + f_{2}(x)\) 几乎处处有定义，且 \(\mu^{*}(f_{1} + f_{2}) \leqslant \mu^{*}(f_{1}) + \mu^{*}(f_{2})\)（化归为在每点 \(f_{1}(x) < +\infty\) 且 \(f_{2}(x) < +\infty\) 的情形，借助 \(a\) )）。
\item
  若 \((f_n)\) 是一个递增的数值函数序列，使得从某个指标起 \(\mu^{*}(f_{n}) > -\infty\)，试证 \(\mu^{*}(\sup_{n} f_{n}) = \sup_{n} \mu^{*}(f_{n})\)。
\end{enumerate}

¶ 6) a) 设 \(f\) 是一个数值函数，使得 \(\mu^{*}(f)\)（习题 5）有限。试证：存在一个可积函数 \(f_{1} \geqslant f\)，使得 \(\mu(f_{1}) = \mu^{*}(f)\)；若 \(f_{2}\) 是另一个可积函数，满足 \(f_{2} \geqslant f\) 且 \(\mu(f_{2}) = \mu^{*}(f)\)，则 \(f_{1}\) 与 \(f_{2}\) 等价。

\begin{enumerate}
\def\labelenumi{\alph{enumi})}
\setcounter{enumi}{1}
\item
  要使数值函数 \(f\) 可积，必须且只须 \(\mu^{*}(f)\) 与 \(\mu_{*}(f)\) 有限且相等。
\item
  设 \(f\) 是一个数值函数，使得 \(\mu^{*}(f)\) 与 \(\mu_{*}(f)\) 都有限；设 \(g\) 与 \(h\) 是两个可积函数，满足 \(g \leqslant f \leqslant h\) 且 \(\mu(g) = \mu_{*}(f)\)、\(\mu(h) = \mu^{*}(f)\)。试证
\end{enumerate}

\[
\mu_ {*} (f - g) = \mu_ {*} (h - f) = 0 \quad \text { and } \quad \mu^ {*} (f - g) = \mu^ {*} (h - f) = \mu^ {*} (f) - \mu_ {*} (f).
\]

\begin{enumerate}
\def\labelenumi{\alph{enumi})}
\setcounter{enumi}{3}
\tightlist
\item
  设 \(f_{1}, f_{2}\) 是两个数值函数，使得数 \(\mu^{*}(f_{1}), \mu^{*}(f_{2})\)、\(\mu_{*}(f_{1})\) 与 \(\mu_{*}(f_{2})\) 都有限；试证
\end{enumerate}

\[
\mu_ {*} (f _ {1} + f _ {2}) \leqslant \mu_ {*} (f _ {1}) + \mu^ {*} (f _ {2}) \leqslant \mu^ {*} (f _ {1} + f _ {2})
\]

（若 \(g_2\) 是一个可积函数，满足 \(f_2 \leqslant g_2\) 且 \(\mu^*(f_2) = \mu(g_2)\)，则注意：对每个可积函数 \(h\)（满足 \(h \leqslant f_1 + f_2\)），有 \(h - g_2 \leqslant f_1\)）。由此推出：若 \(f_1\) 可积，则

\[
\mu^ {*} (f _ {1} + f _ {2}) = \mu (f _ {1}) + \mu^ {*} (f _ {2}), \quad \mu_ {*} (f _ {1} + f _ {2}) = \mu (f _ {1}) + \mu_ {*} (f _ {2}),
\]

且 \(\mu^{*}(f_{1} + f_{2}) = \mu^{*}\big(\sup (f_{1},f_{2})\big) + \mu^{*}\big(\inf (f_{1},f_{2})\big)\)（对后一关系式，利用 §1 的习题 1）。

\begin{enumerate}
\def\labelenumi{\alph{enumi})}
\setcounter{enumi}{4}
\tightlist
\item
  设 \(f\) 是一个可积函数。要使函数 \(g\)（\(\mu^{*}(g)\) 有限）可积，必须且只须 \(\mu(f) = \mu^{*}(g) + \mu^{*}(f - g)\)（若 \(g_{1}\) 是一个可积函数，满足 \(g \leqslant g_{1}\) 且 \(\mu^{*}(g) = \mu(g_{1})\)，注意 \(f - g_{1} \leqslant f - g\)）。
\end{enumerate}

¶ 7) 设 \(\mu\) 是 X 上的一个正测度。对每个子集 \(A \subset X\)，特征函数 \(\varphi_A\) 的下积分（习题 5）称为 A 的内测度，记作 \(\mu_*(A)\)。

\begin{enumerate}
\def\labelenumi{\alph{enumi})}
\item
  试证：\(\mu_{*}(\mathrm{A})\) 是含于 A 的紧集的测度之上确界（按定理 4 的方式论证）。
\item
  对 X 的每个外测度有限的子集 A，试证：存在两个可积集 \(A_{1}, A_{2}\)，使得 \(A_{1} \subset A \subset A_{2}\) 且 \(\mu_{*}(A) = \mu(A_{1})\)、\(\mu^{*}(A) = \mu(A_{2})\)。要使 A 可积，必须且只须 \(\mu^{*}(A)\) 与 \(\mu_{*}(A)\) 有限且相等。用同样的记号，试证
\end{enumerate}

\[
\mu_ {*} (\mathrm{A} \cap \mathbb {C} \mathrm{A} _ {1}) = \mu_ {*} (\mathrm{A} _ {2} \cap \mathbb {C} \mathrm{A}) = 0
\]

以及

\[
\mu^ {*} (\mathrm{A} \cap \mathbb {C} \mathrm{A} _ {1}) = \mu^ {*} (\mathrm{A} _ {2} \cap \mathbb {C} \mathrm{A}) = \mu^ {*} (\mathrm{A}) - \mu_ {*} (\mathrm{A}).
\]

\begin{enumerate}
\def\labelenumi{\alph{enumi})}
\setcounter{enumi}{2}
\item
  设 A 是一个可积集；试证：对每个集 \(\mathbf{B} \subset \mathbf{A}\)，有 \(\mu(\mathbf{A}) = \mu^{*}(\mathbf{B}) + \mu_{*}(\mathbf{A} \cap \mathbb{C}\mathbf{B})\)。
\item
  设 A 与 B 是两个不相交的外测度有限的集。试证：若 \(\mathbf{C} = \mathbf{A} \cup \mathbf{B}\)，则
\end{enumerate}

\[
\mu_ {*} (\mathrm{A}) + \mu_ {*} (\mathrm{B}) \leqslant \mu_ {*} (\mathrm{C}) \leqslant \mu_ {*} (\mathrm{A}) + \mu^ {*} (\mathrm{B}) \leqslant \mu^ {*} (\mathrm{C}) \leqslant \mu^ {*} (\mathrm{A}) + \mu^ {*} (\mathrm{B})
\]

以及

\[
\mu_ {*} (\mathrm{C}) - \mu_ {*} (\mathrm{A}) - \mu_ {*} (\mathrm{B}) \leqslant \mu^ {*} (\mathrm{A}) + \mu^ {*} (\mathrm{B}) - \mu^ {*} (\mathrm{C})
\]

（对后一不等式，化归为 \(\mu_{*}(\mathrm{A}) = \mu_{*}(\mathrm{B}) = 0\) 的情形，借助 \(b\) )）；若 \(\mathrm{A}_2\) 与 \(\mathrm{B}_2\) 是可积集，满足 \(\mathrm{A} \subset \mathrm{A}_2\)、\(\mathrm{B} \subset \mathrm{B}_2\)、\(\mu^{*}(\mathrm{A}) = \mu(\mathrm{A}_2)\)、\(\mu^{*}(\mathrm{B}) = \mu(\mathrm{B}_2)\)，试证 \(\mu_{*}(\mathrm{C}) \leqslant \mu(\mathrm{A}_2 \cap \mathrm{B}_2)\)。

¶ 8) 若把环面 T = R/Z 与 R 的区间 {[}0,1 {[} 典范地等同起来，该区间上的 Lebesgue 测度 μ 就是 T 上的一个测度；对每个集 A ⊂ T 及每个 z ∈ T，有 μ*(A + z) = μ*(A)。

\begin{enumerate}
\def\labelenumi{\alph{enumi})}
\item
  试证：存在子群 \(\mathrm{H}_0\)（\(\mathbf{T}\) 的），使得诸集 \(\mathrm{H}_n = r_n + \mathrm{H}_0\)（其中 \(r_n\) 取遍含于 {[}0,1{[} 的有理数所成之集）构成 \(\mathbf{T}\) 的一个分割（把 \(\mathbf{R}\) 看作 \(\mathbf{Q}\) 上的向量空间，注意在 \(\mathbf{R}\) 中子空间 \(\mathbf{Q}\) 有一个补子空间）。
\item
  设 H 是任意有限多个集 \(H_{n}\) 的并组成的集；试证：H 不可积且 \(\mu_{*}(H)=0\)（注意加群 Q/Z 中由有限个元素生成的每个子群在 Q/Z 中都有无限指数；由此推出，存在 T 的一个分成可数无穷多个集 \(P_{n}\) 的分割，其中每个集都含有形如 \(z+H\) 的集）。
\item
  由 \(b\) ) 出发，给出一个递减序列 \((\mathrm{A}_n)\)（\(\mathbf{T}\) 的子集的序列）的例子，其交为空集，且 \(\mu^{*}(\mathrm{A}_{n}) = 1\)（对所有 \(n\)）。
\item
  设 A 是一个可积集，满足 \(H_{0} \subset A\) 且 \(\mu(A) = \mu^{*}(H_{0}) > 0\)；试证：对每个可积集 \(B \subset A\)，集 \(B_{1} = B \cap H_{0}\) 与 \(B_{2} = B \cap C H_{0}\) 构成 B 的一个分割，使得 \(\mu^{*}(B_{1}) = \mu^{*}(B_{2}) = \mu(B)\) 且 \(\mu_{*}(B_{1}) = \mu_{*}(B_{2}) = 0\)（利用习题 7 b)）。
\end{enumerate}

¶ 9) a) 设 \(\mu\) 是局部紧空间 \(X\) 上的一个正测度。在 Banach 空间 \(\widetilde{\mathcal{F}}^1\)（由 \(\mathcal{F}^1\) 中函数的等价类组成）中，设 \(G\) 是一个包含子空间 \(L^{1}\) 的闭线性子空间；设 G 是由类属于 G 的函数组成的 \(F^{1}\) 的闭线性子空间。试证：可把线性形式 \(\mu(f)\) 延拓到 G 上，使不等式 \(|\mu(f)| \leqslant N_{1}(f)\) 仍成立（利用 Hahn--Banach 定理）。试证：对这样延拓的积分，§3 的定理 3 仍然有效。

\begin{enumerate}
\def\labelenumi{\alph{enumi})}
\setcounter{enumi}{1}
\tightlist
\item
  假设 G 是 \(\mathbf{L}^1\) 与一个有限维子空间 H 的直和。试证：在 \(\mathcal{G}\) 中，Lebesgue 定理（定理 2）仍然有效。（设 \((f_n)\) 是 \(\mathcal{G}\) 中函数的一个序列，几乎处处趋于 \(f\)，且满足 \(|f_n| \leqslant g\)，其中 \(g \geqslant 0\) 且 \(\mathrm{N}_1(g) < +\infty\)。设 \((\widetilde{u}_k)_{1 \leqslant k \leqslant m}\) 是 H 的一个基，并设 \(f_n = \sum_{k=1}^{m} \alpha_{nk} u_k + h_n\)，其中 \(h_n \in \mathcal{L}^1\)。利用 \(G\) 是 \(L^1\) 与 \(H\) 的拓扑直和这一事实，试证诸 \(\alpha_{nk}\) 一致有界；必要时过渡到 \((f_n)\) 的一个子序列，化归为每个序列 \((\alpha_{nk})_{n \geqslant 1}\) 都有极限的情形；由此推出 \(h_n(x)\) 这时几乎处处趋于一个极限，并对序列 \((h_n)\) 应用 Lebesgue 定理；最后注意 \((\widetilde{f}_n)\) 的两个子序列在 \(G\) 中不可能趋于不同的极限。
\end{enumerate}

¶ 10) a) 设 E 是一个 Riesz 空间，μ 是 E 上的一个正线性形式，使得关系 μ(\textbar x\textbar) = 0 蕴含 x = 0；于是 μ(\textbar x\textbar) 是 E 上的一个范数，我们假设 E 赋此范数后是完备的。由此推出 E 是完全格序的（Ch. II, §2, 习题 8 e)）。从而（Ch. II, §1, 习题 12 与 13）存在一个局部紧空间 X，它是一族 (Kα) 紧 Stone 空间之和，使得 E 同构于 X 上的一个包含空间 K(X) 的连续数值函数（无论有限与否）所成的空间。把 E 与这个空间等同起来，μ 在 K(X) 上的限制就是 X 上的一个正测度。试证：E 典范地同构于空间 L¹(μ)；更精确地说，对每个定义在 X 上的 μ-可积函数 g，存在一个且仅一个函数 f ∈ E，它关于测度 μ 与 g 等价（注意 E 的每个元素 ≥ 0 都是 K(X) 的元素的一个递增序列的上确界，且 K(X) 在 L¹(μ) 中稠密）。

\begin{enumerate}
\def\labelenumi{\alph{enumi})}
\setcounter{enumi}{1}
\tightlist
\item
  由 a) 推出：对每个紧空间 K，存在一个局部紧空间 S（一族紧 Stone 空间的拓扑和）及 S 上的一个正测度 \(\nu\)（其支集等于 S），使得 K 上的测度所成的完全格序空间 \(\mathcal{M}(\mathrm{K})\)（赋以范数 \(\|\mu\|\)）同构于 \(\mathcal{L}^{1}(\nu)\)（考虑线性形式 \(\mu\mapsto\mu(\mathrm{K})\)，在 \(\mathcal{M}(\mathrm{K})\) 上）。
\end{enumerate}

\begin{enumerate}
\def\labelenumi{\arabic{enumi})}
\setcounter{enumi}{10}
\tightlist
\item
  \begin{enumerate}
  \def\labelenumii{\alph{enumii})}
  \tightlist
  \item
    设 \(\Gamma\) 是一个集 \(A\) 的子集的任意一个集合。设 \(\Psi\) 是 \(A\) 的形如下式的子集所成之集
  \end{enumerate}
\end{enumerate}

\[
\mathrm{X} _ {1} \cap \mathrm{X} _ {2} \cap \dots \cap \mathrm{X} _ {m} \cap \mathbb {C} \mathrm{X} _ {m + 1} \cap \dots \cap \mathbb {C} \mathrm{X} _ {m + p},
\]

其中诸 \(X_{i}\) 是 \(\Gamma\) 中的集，m 是任意整数 \(\geqslant 1\)，p 是任意整数 \(\geqslant 0\)。试证：最小集环 \(\Phi\)（包含 \(\Gamma\)）是 \(\Psi\) 中之集的有限并所成之集。

\begin{enumerate}
\def\labelenumi{\alph{enumi})}
\setcounter{enumi}{1}
\item
  设 \(\Delta\) 是 \(\Gamma\) 中之集的有限交所成之集。试证：对每个向量空间 F，由 \(\Gamma\) 生成的集环中之集的特征函数的线性组合（系数取自 F）所成之集，等同于 \(\Delta\) 中之集的特征函数的线性组合所成之集。
\item
  设 X 是拓扑空间，\(\Gamma\) 是 X 的紧子集所成之集。试证：由 \(\Gamma\) 生成的集环等同于 X 的形如 \(A \cap CB\) 的子集的有限并所成之集，其中 A 与 B 是紧集。
\end{enumerate}

\begin{enumerate}
\def\labelenumi{\arabic{enumi})}
\setcounter{enumi}{11}
\tightlist
\item
  设 \(X\) 是一个 Hausdorff 拓扑空间。对每个偶 \((K, U)\)（由一个紧集 \(K\) 与一个开集 \(U\)——\(X\) 中的——组成），设 \(I(K, U)\) 是子集 \(M \subset X\)（满足 \(K \subset M \subset U\)）所成之集，并设 \(\mathcal{T}\) 是由 \(\mathfrak{P}(X)\) 的诸子集 \(I(K, U)\) 在 \(\mathfrak{P}(X)\) 上生成的拓扑。
\end{enumerate}

\begin{enumerate}
\def\labelenumi{\alph{enumi})}
\item
  试证：每个集 I(K, U) 在 \(\mathfrak{P}(X)\) 中既开又闭；由此推出，\(\mathfrak{P}(X)\) 赋以拓扑 \(\mathcal{T}\) 后是一个完全正则且全不连通的空间。
\item
  要使映射 \(M \mapsto CM\)（从 \(\mathfrak{P}(X)\) 到其自身）连续（关于拓扑 T），必须且只须 X 是紧的。
\item
  取 X 为 R 的区间 {[}0,1{]}。试证：映射 \((M,N)\mapsto M\cup N\) 与 \((M,N)\mapsto M\cap N\)（从 \(\mathfrak{P}(X)\times\mathfrak{P}(X)\) 到 \(\mathfrak{P}(X)\) 中）关于拓扑 T 不连续。
\item
  假设 X 是局部紧的。试证：T 在 X 的紧子集所成之集上诱导的拓扑，细于按 GT, II, §1 习题 5 的程序由 X 上一个一致结构导出的拓扑；除非 X 是离散空间，这两个拓扑不可能相同。
\end{enumerate}

¶ 13) a) 设 X 是局部紧空间，\(\Gamma\) 是由相对紧的集组成的 X 拓扑的一个基。设 \(\mathcal{U}\) 是与 X 的拓扑相容的一个一致结构，\(\mathfrak{S}\) 是该结构的一个邻域基本系。设 \(M \mapsto \lambda(M)\) 是一个有限数值函数 \(\geqslant 0\)（定义在 \(\Gamma\) 上）。对每个紧集 \(K \subset X\) 与每个邻域 \(V \in \mathfrak{S}\)，设 \(\alpha_V(K)\) 是数 \(\sum_{i} \lambda(U_i)\)

对 K 的所有有限覆盖 \((\mathrm{U}_{i})\)（由 \(\Gamma\) 中关于 V 为细小的集组成）；假设当 V 取遍 S 时，\(\alpha(\mathrm{K})\)（诸数 \(\alpha_{\mathrm{V}}(\mathrm{K})\) 的上确界）有限。试证：存在 X 上的一个测度 \(\mu\)（且仅一个），使得对每个紧集 K 有 \(\mu(\mathrm{K}) = \alpha(\mathrm{K})\)（利用定理 5）。

\begin{enumerate}
\def\labelenumi{\alph{enumi})}
\setcounter{enumi}{1}
\tightlist
\item
  设 \(\Psi\) 是 \(X\) 的 Borel 子集所成之集，\(\alpha\) 是 \(\Psi\) 到 \([0, +\infty]\) 中的一个映射，满足下列条件：
\end{enumerate}

\begin{enumerate}
\def\labelenumi{(\roman{enumi})}
\item
  若 \(\mathrm{B}_1, \mathrm{B}_2\) 是 \(X\) 的两个不交的 Borel 子集，则 \(\alpha(\mathrm{B}_1 \cup \mathrm{B}_2) = \alpha(\mathrm{B}_1) + \alpha(\mathrm{B}_2)\)。\\
\item
  若 \(B\) 是 \(X\) 的一个紧子集，则 \(\alpha(B) < +\infty\)。
\item
  若 \(\mathbf{B}\) 是 \(\mathbf{X}\) 的一个 Borel 子集，则 \(\alpha(\mathbf{B}) = \inf \alpha(\mathbf{U})\)，其中 \(\mathbf{U}\) 取遍 \(\mathbf{X}\) 中包含 \(\mathbf{B}\) 的开子集所成之集。
\end{enumerate}

那么，存在一个且仅一个 X 上的正测度 \(\mu\)，使得对 X 的每个可被一列紧集覆盖的 Borel 子集 B 有 \(\alpha(B) = \mu^{*}(B)\)。c) 对每个 \(B \in \Phi\)，若 B 可被一列紧集覆盖则令 \(\alpha(B) = 0\)，在相反情形令 \(\alpha(B) = +\infty\)。这时 \(b\) 的条件 (i)、(ii)、(iii) 满足。我们有 \(\mu = 0\)，因此若 B 不能被一列紧集覆盖，则 \(\alpha(B) \neq \mu^{*}(B)\)。

\begin{enumerate}
\def\labelenumi{\arabic{enumi})}
\setcounter{enumi}{13}
\item
  设 \((\mathbf{A}_n)\) 是一个可积集序列，使得 \(\sum_{n} \mu(\mathbf{A}_n) < +\infty\)。对每个整数 \(k\)，设 \(\mathbf{G}_k\) 是由 \(x \in X\) 组成的集，其中 \(x \in \mathbf{A}_n\) 至少对 \(k\) 个值 \(n\) 成立；试证：\(\mathbf{G}_k\) 可积，且 \(k \cdot \mu(\mathbf{G}_k) \leqslant \sum_{n} \mu(\mathbf{A}_n)\)。
\item
  设 \(\mu\) 是局部紧空间 \(X\) 上的一个有界正测度，\((A_n)\) 是 \(X\) 中可积集的一个序列，使得 \(\inf_{n} \mu(A_n) = m > 0\)。试证：集 \(B\)——\(X\) 中属于无穷多个集 \(A_n\) 的点所成之集——是可积的，且 \(\mu(B) \geqslant m\)。
\end{enumerate}

¶ 16) 设 X 是一个完全正则空间，\(\mathcal{C}(X)\)（相应地，\(\mathcal{C}^b(X)\)）是 X 上连续（相应地，连续且有界）数值函数的 Riesz 空间。

\begin{enumerate}
\def\labelenumi{\alph{enumi})}
\item
  试证：若线性形式 \(\lambda\)（\(\mathcal{C}(\mathrm{X})\) 上的）关于紧收敛拓扑连续，则它是相对有界的。
\item
  设 \(\lambda\) 是 \(\mathcal{C}(\mathrm{X})\) 上的一个正线性形式。试证：若 \(\lambda\) 在 \(\mathcal{C}^b (\mathrm{X})\) 上为零，则它在 \(\mathcal{C}(\mathrm{X})\) 上为零（设 \(\varphi\) 是 \(\mathcal{C}(\mathrm{X})\) 到商 Riesz 空间 \(\mathcal{C}(\mathrm{X}) / \mathcal{C}^b (\mathrm{X})\)（Ch. II, §1, 习题 4）上的典范映射；试证：若 \(h\) 是一个连续函数 \(\geqslant 0\) 且在 \(\mathrm{X}\) 上不有界，则 \(n\varphi (h)\leqslant \varphi (h^{2})\)（对每个整数 \(n > 0\)））。
\item
  设 \(\beta X\) 是 \(X\) 的 Stone--Čech 紧化，即通过在 \(X\) 上赋予以使 \(\mathcal{C}^b(X)\) 中函数都一致连续的最粗一致结构、再把如此得到的一致空间完备化而得的紧空间；这时每个函数 \(f \in \mathcal{C}(X)\) 都可连续延拓为连续函数 \(\widetilde{f}\)（在 \(\beta X\) 上，可能取无穷值）（考虑 \(f/(1 + |f|)\)）。若 \(\lambda\) 是 \(\mathcal{C}(X)\) 上的一个正线性形式，则 \(\lambda\) 在 \(\mathcal{C}^b (\mathrm{X})\) 上的限制形如 \(f\mapsto \mu (\widetilde{f})\)，其中 \(\mu\) 是 \(\beta X\) 上的一个正测度；试证：对每个函数 \(f\in \mathcal{C}(\mathrm{X})\)，\(\widetilde{f}\) 关于 \(\mu\) 可积，且 \(\lambda (f) = \mu (\widetilde{f})\)（利用 \(b\) )，并注意每个函数 \(\geqslant 0\)（属于 \(\mathcal{C}(\mathrm{X})\) 的）都是 \(\mathcal{C}^b (\mathrm{X})\) 中函数的一个序列的上包络）。
\item
  试证：每个点 \(x_{0}\)（属于 \(\mu\) 的支集而不属于 X 的）都具有下列性质：对每个递减序列 \((V_{n})\)（由 \(x_{0}\) 的邻域组成），诸 \(V_{n}\) 的交至少含有 X 的一个点。逆命题成立。
\item
  由 \(d\) ) 推出：若 \(X\) 局部紧且无穷远处可数，则 \(\mu\) 的支集含于 \(X\)（与 \(h\) ) 比较）。
\item
  若 X 是离散的，试证：\(\mu\) 的支集有限（在相反情形，造一个定义在 X 上的函数 \(f \geqslant 0\)，使 \(\widetilde{f}\) 不是 \(\mu\) -可积的）。
\item
  试证：若线性形式 \(\lambda\)（\(\mathcal{C}(\mathbf{X})\) 上的）是正的且关于紧收敛拓扑连续，则 \(\mu\) 的支集含于 X（参见 Ch. III, §2, No.~3, 命题 11）。
\item
  设 \(X_{0}\) 是把 GT, I, §9, 习题 11 中定义的局部紧空间 X 添加一个无穷远点 \(\omega\) 所得的紧空间，Y 是 \(\overline{R}\) 中由整数 \(\geqslant 0\) 与 \(+\infty\) 组成的子空间，Z 是点 \((\omega, +\infty)\) 在乘积空间 \(X_{0} \times Y\) 中的余空间（局部紧）。试证：Z 上每个连续数值函数都有界，且 \(f(z)\) 当 z 趋于 Z 的无穷远点 \((\omega, +\infty)\) 时趋于一个极限。由此推出：若 \(\mu\) 是 Z 上由把质量 \(1/2^{n}\) 放在每个点 \((\omega, n)\) 处（对每个 \(n \geqslant 0\)）所定义的测度，则 Z 上每个连续函数都 \(\mu\) -可积，但 \(\mu\) 的支集不是紧的。
\end{enumerate}

¶ 17) 设 X 是 GT, I, §9, 习题 11 中定义的局部紧空间。

\begin{enumerate}
\def\labelenumi{\alph{enumi})}
\tightlist
\item
  设 H 是赋以紧收敛拓扑的局部凸空间 \(\mathcal{C}(\mathrm{X};\mathbf{C})\) 的一个紧子集；试证：诸函数 \(f \in H\) 在 X 上一致有界，且存在 \(c \in X\)，使得对每个 \(x \geqslant c\)，H 中所有函数都取常值。（用反证法论证；若 H 中函数在 X 上不一致有界，就会存在 X 的点的一个递增序列 \((x_{n})\)，使得 \(\sup(|\mathrm{H}(x_{n})|) \geqslant n\)，其中 \(|\mathrm{H}(x_{n})|\) 表示诸数 \(|f(x_{n})|\)（对 \(f \in H\)）所成之集；注意序列 \((x_{n})\) 在 X 中收敛。类似地，对每个 \(x \in X\)，设 \(\delta(x)\) 是所有函数 \(f \in H\) 在区间 \([x, \rightarrow[; \delta(x)\) 上的振幅的上确界；它有限且递减，因此存在 \(d \in X\)，使得 \(\delta(x)\) 对 \(x \geqslant d\) 取常值。为证这个常数 \(\beta\) 为零，像前面那样用反证法，造出 X 的点的两个递增序列 \((s_{n}), (t_{n})\)，满足 \(s_{n} \leqslant t_{n} \leqslant s_{n+1} \leqslant t_{n+1}\)，以及 H 中函数的一个序列 \((f_{n})\)，满足 \(|f_{n}(s_{n}) - f_{n}(t_{n})| \geqslant \beta/2\)。
\end{enumerate}

特别地，\(\mathcal{C}(\mathrm{X};\mathbf{C})\) 是 \(\mathcal{K}(\mathrm{X};\mathbf{C})\) 与 \(C\cdot1\) 的直和。

\begin{enumerate}
\def\labelenumi{\alph{enumi})}
\setcounter{enumi}{1}
\item
  试证：在 X 上每个测度都有紧支集（对每个 \(x \in X\)，设 \(f_x\) 是区间 \(\leftarrow, x\) {]} 的特征函数，它连续且具有紧支集；考虑 X 上的递增函数 \(|\mu|(f_x)\)）。
\item
  由 a) 与 b) 推出：在 \(\mathcal{M}(\mathrm{X};\mathbf{C}) = \mathcal{M}^{1}(\mathrm{X};\mathbf{C}) = \mathcal{C}'(\mathrm{X};\mathbf{C})\) 中，关于 \(\mathcal{T}\)（\(\mathcal{C}(\mathrm{X};\mathbf{C})\) 的紧子集上一致收敛拓扑）的有界集，就是关于超强拓扑（Ch. III, §1, 习题 15）的有界集。试证：\(\mathcal{C}'(\mathrm{X};\mathbf{C})\) 关于拓扑 \(\mathcal{T}\) 不是拟完备的（考虑测度 \(\varepsilon_{x}\)，对 \(x\in \mathrm{X}\)）；由此推出，关于紧收敛拓扑，\(\mathcal{C}(\mathrm{X};\mathbf{C})\) 既非桶式也非有界型（参见 TVS, III, §4, No.~2, 定理 1 的推论 4 及下面的习题 18）。
\item
  试证：在 \(\mathcal{K}(\mathrm{X};\mathbf{C})\) 上，拓扑 \(\mathcal{T}_0\)（诸 Banach 空间 \(\mathcal{K}(\mathrm{X},\mathrm{K};\mathbf{C})\) 的拓扑的直接极限）等同于一致收敛拓扑，且赋此拓扑后 \(\mathcal{K}(\mathrm{X};\mathbf{C})\) 是完备的。（设 V 是 \(\mathcal{T}_0\) 的一个 0 的邻域；对每个 \(x\in X\)，设 \(r_x\) 是最大的数 \(>0\)，使得 V 包含每个连续函数 \(f\)，其支集含于 \([ \leftarrow ,x]\) 且满足 \(\| f\| \leqslant r_x\)。试证 \(r_x\) 在 X 中的下确界 \(>0\)，为此证明不可能存在 X 的点的递增序列 \((x_n)\)，使得 \(\lim_{n\to \infty}r_{x_n} = 0\)。）
\end{enumerate}

\begin{enumerate}
\def\labelenumi{\arabic{enumi})}
\setcounter{enumi}{17}
\tightlist
\item
  设 \(X\) 是一个仿紧的局部紧空间。
\end{enumerate}

\begin{enumerate}
\def\labelenumi{\alph{enumi})}
\item
  试证：空间 \(\mathcal{C}(\mathrm{X};\mathbf{C})\) 同构于一族 Fréchet 空间的乘积，从而是桶式的。
\item
  要使 \(\mathcal{C}^{\prime}(X; \mathbf{C})\) 的子集 H 关于紧收敛拓扑有界，必须且只须存在 X 的一个紧子集 K，使得 \(\operatorname{Supp}(\mu) \subset K\)（对一切 \(\mu \in H\)），且 H 关于超强拓扑有界（与习题 17 c) 比较）。这时紧收敛拓扑在 H 上诱导的拓扑等同于淡拓扑。
\item
  试证：在集 \(\mathcal{M}_{+}(\mathrm{X})\cap \mathcal{C}'(\mathrm{X};\mathbf{C})\) 上，由 \(\mathcal{C}'(\mathrm{X};\mathbf{C})\) 上紧收敛拓扑诱导的拓扑等同于淡拓扑。当 \(X\) 不是仿紧时（习题 17），结论还成立吗？
\end{enumerate}

\begin{enumerate}
\def\labelenumi{\arabic{enumi})}
\setcounter{enumi}{18}
\tightlist
\item
  在 \(\mathcal{C}(\mathrm{X};\mathbf{C})\) 上赋以紧收敛拓扑，在 \(\mathcal{C}'(\mathrm{X};\mathbf{C})\) 上赋以 \(\mathcal{C}(\mathrm{X};\mathbf{C})\) 的紧子集上一致收敛拓扑。设 \(\mathfrak{S}\) 是 \(\mathcal{C}(\mathrm{X};\mathbf{C})\) 的紧子集所成之集；试证：双线性映射 \((g,\mu)\mapsto g\cdot \mu\)（从 \(\mathcal{C}(\mathrm{X};\mathbf{C})\times \mathcal{C}'(\mathrm{X};\mathbf{C})\) 到 \(\mathcal{C}'(\mathrm{X};\mathbf{C})\) 中）是 \(\mathfrak{S}\) -亚连续的。设 \(\mathcal{T}\) 是 \(\mathcal{C}'(\mathrm{X};\mathbf{C})\) 的有界子集所成之集；试证：若 X 是仿紧的，则前面的双线性映射也是 \(\mathcal{T}\) -亚连续的；当不再假设 X 是仿紧时（习题 17），后一性质还成立吗？
\end{enumerate}

¶ 20) 设 X, Y 是两个仿紧的局部紧空间。在空间 \(\mathcal{C}'(X; \mathbf{C})\) 与 \(\mathcal{C}'(Y; \mathbf{C})\) 上赋以紧收敛拓扑。

\begin{enumerate}
\def\labelenumi{\alph{enumi})}
\item
  试证：当在 \(\mathcal{C}^{\prime}(X\times Y;\mathbf{C})\) 上赋以紧收敛拓扑时，映射 \((\mu,\nu)\mapsto\mu\otimes\nu\) 是 \((\mathfrak{S},\mathfrak{T})\) -亚连续的，其中 S 表示 \(\mathcal{C}^{\prime}(X;\mathbf{C})\) 的有界子集所成之集，T 表示 \(\mathcal{C}^{\prime}(Y;\mathbf{C})\) 的有界子集所成之集（利用习题 18 b)）。
\item
  试证：当 X 与 Y 无穷远处可数且在 \(\mathcal{C}^{\prime}(X\times Y;C)\) 上赋以紧收敛拓扑时，映射 \((\mu,\nu)\mapsto\mu\otimes\nu\) 连续。（受 Ch. III, §4 习题 3 的证明启发，在 X（相应地 Y）中引入从属于一个局部有限开覆盖的连续单位分解，按 TVS, II, §4, 习题 9 a) 的方式论证。）
\item
  若 X 是离散的，空间 \(\mathcal{C}^{\prime}(X;\mathbf{C})\) 可等同于直和空间 \(\mathbf{C}^{(X)}\)（赋以最细局部凸拓扑）。由此给出一个例子：X 与 Y 都是离散的，X 不可数，且映射 \((\mu,\nu)\mapsto\mu\otimes\nu\) 当 \(\mathcal{C}^{\prime}(X\times Y;\mathbf{C})\) 赋以淡拓扑时不连续（参见 TVS, II, §4, 习题 9 b)）。
\end{enumerate}

\begin{enumerate}
\def\labelenumi{\arabic{enumi})}
\setcounter{enumi}{20}
\item
  设 \(\mathbf{X}\) 是一个局部紧空间，\(f\) 是一个数值函数 \(\geqslant 0\)（定义在 \(\mathbf{X}\) 上）。要使 \(\mu \mapsto \mu^{*}(f)\) 从 \(\mathcal{M}_{+}(\mathbf{X}) \cap \mathcal{C}'(\mathbf{X};\mathbf{C})\) 到 \(\overline{\mathbf{R}}\) 中的映射关于紧收敛拓扑连续，必须且只须 \(f\) 在 \(\mathbf{X}\) 上连续。
\item
  设 X, Y 是两个局部紧空间，\(\mu\) 是 X 上的一个有界测度，\(\nu\) 是 Y 上的一个有界测度，于是 \(\mu \otimes \nu\) 是 \(X \times Y\) 上的一个有界测度。试证：对每个函数 \(f \in \mathcal{C}^{0}(X \times Y; \mathbf{C})\)，函数 \(x \mapsto \int f(x, y) d\nu(y)\) 属于 \(\mathcal{C}^{0}(X; \mathbf{C})\)，并且
\end{enumerate}

\[
\int f (x, y) d \mu (x) d \nu (y) = \int d \mu (x) \int f (x, y) d \nu (y).
\]

¶ 23) 在空间 \(\mathbf{R}^n\) 上，设 \(\mu\) 是 Lebesgue 测度，\(|\mathbf{x}|\) 是一个范数，使该范数的单位球测度等于 1，\((\mathbf{x}_k)_{k \geqslant 1}\) 是一个有界可积集 B 中互异点的一个无穷序列，其中 \(\mu(B) = 1\)。对每个整数 \(m\)，用 \(d_m\) 表示数 \(|\mathbf{x}_i - \mathbf{x}_j|\)（对 \(1 \leqslant i < j \leqslant m\)）的下确界。试证

\(\liminf_{m \to \infty} md_m^n \leqslant \alpha_n^{-1}\)，其中

\[
\alpha_ {n} = 1 + n \int_ {0} ^ {1} \frac {(1 - t) ^ {n}}{1 + t} d t.
\]

（用反证法论证：假设对某个 \(\varepsilon > 0\)，存在一个 \(m_0\)，使得 \(md_m^n > h_n\)（对 \(m \geqslant m_0\)），其中 \(h_n = \alpha_n^{-1} + \varepsilon\)。对 \(1 \leqslant i < m\)，设 \(B_i\) 是以 \(\mathbf{x}_i\) 为中心、以 \(\frac{1}{2} h_n m^{-1/n}\) 为半径的球；对 \(m \leqslant i \leqslant 2^n m\)，设 \(B_i\) 是以 \(\mathbf{x}_i\) 为中心、以 \(\frac{1}{2} h_n (2i^{-1/n} - m^{-1/n})\) 为半径的球。试证：这 \(2^n m\) 个球 \(B_i\) 两两不交，并利用 Euler--Maclaurin 公式计算它们的并的测度。）
